\documentclass[11pt,a4paper]{article}

\usepackage[utf8]{inputenc}
\usepackage[T1]{fontenc}
\usepackage{lmodern}
\usepackage[french]{babel}
\usepackage[margin=2.5cm]{geometry}
\usepackage{microtype}          % justification fine, moins de débordements
\usepackage{booktabs}
\usepackage{array}
\usepackage[table]{xcolor}
\usepackage{colortbl}
\usepackage{amsmath}
\usepackage{amssymb}
\usepackage{enumitem}
\usepackage{titlesec}
\usepackage{fancyhdr}
\usepackage{graphicx}
\usepackage{caption}
\usepackage[official]{eurosym}
\usepackage[round,authoryear]{natbib}
\usepackage{hyperref}           % chargé en dernier (bonne pratique)

\definecolor{navy}{HTML}{1F3864}
\definecolor{bluemid}{HTML}{2E5496}
\definecolor{lightblue}{HTML}{D6E0F0}

\hypersetup{
  colorlinks=true, linkcolor=bluemid, urlcolor=bluemid, citecolor=bluemid,
  bookmarksnumbered=true, bookmarksopen=true,
  pdftitle={Quantification du SCR lié à la non-conformité au règlement DORA},
  pdfauthor={Kélian Kaddouri}
}

\captionsetup{font=small,labelfont={bf,color=navy},labelsep=period}

\titleformat{\section}{\color{navy}\large\bfseries}{\thesection}{0.6em}{}
\titleformat{\subsection}{\color{bluemid}\normalsize\bfseries}{\thesubsection}{0.5em}{}
\titleformat{\subsubsection}{\normalsize\bfseries}{\thesubsubsection}{0.5em}{}
% Respiration verticale autour des titres
\titlespacing*{\section}{0pt}{2.4ex plus 1ex minus .2ex}{1.4ex plus .2ex}
\titlespacing*{\subsection}{0pt}{2.0ex plus .8ex minus .2ex}{1.0ex plus .2ex}

\pagestyle{fancy}
\fancyhf{}
\lhead{\small\color{navy}Mémoire SCR\_DORA}
\chead{\small\color{navy}Quantification du SCR lié à DORA}
\rhead{\small\thepage}
\renewcommand{\headrulewidth}{0.4pt}

\setlength{\parskip}{0.4em}
\setlength{\parindent}{1.2em}
\setlength{\emergencystretch}{3em}   % réduit les débordements résiduels
% Titres de figures/tables plus aérés
\setlength{\intextsep}{14pt plus 2pt minus 2pt}
\captionsetup{skip=6pt}

\begin{document}

%==============================================================
% PAGE DE GARDE
%==============================================================
\begin{titlepage}
\pagestyle{empty}
\begin{center}

\vspace*{1cm}
{\small\color{navy}\textsc{ENSAE Paris}}\\
{\small\color{navy}\textsc{Nexialog Consulting}}

\vspace{2cm}
\noindent\rule{\textwidth}{0.8pt}\\[0.5cm]

{\LARGE\color{navy}\bfseries Quantification du SCR lié à la\\[4pt] non-conformité au règlement DORA}\\[10pt]
{\large\color{bluemid} Un premier pont quantitatif entre résilience opérationnelle numérique et solvabilité prudentielle}

\vspace{0.5cm}
\noindent\rule{\textwidth}{0.8pt}

\vspace{2.5cm}
{\large Mémoire d'actuariat}

\vspace{2cm}

\begin{tabular}{rl}
\textbf{Auteur :} & Kélian Kaddouri \\[4pt]
\textbf{Établissement :} & ENSAE Paris --- Master Actuariat \\[4pt]
\textbf{Structure d'accueil :} & Nexialog Consulting \\[4pt]
\textbf{Encadrant professionnel :} & Hugo Rapior \\[4pt]
\textbf{Référente académique :} & Caroline Hillairet \\
\end{tabular}

\vfill

{\small\color{navy} Version de travail --- \today}

\end{center}
\end{titlepage}

%==============================================================
% TABLE DES MATIÈRES
%==============================================================
\begingroup
\hypersetup{linkcolor=navy}
\renewcommand{\contentsname}{\color{navy}Sommaire}
\setcounter{tocdepth}{2}   % sections + sous-sections seulement (sommaire épuré)
\tableofcontents
\endgroup
\thispagestyle{fancy}
\clearpage

\setcounter{page}{1}

%==============================================================
\section*{Introduction}
\addcontentsline{toc}{section}{Introduction}
%==============================================================

Le risque cyber occupe aujourd'hui une place singulière dans le paysage
prudentiel des institutions financières. Sa nature hybride --- à la fois
opérationnelle, technologique et parfois systémique --- le place à la
croisée de plusieurs cadres réglementaires qui, historiquement, ne
dialoguaient pas directement entre eux. En assurance, le risque opérationnel
demeure traité dans la formule standard de Solvabilité~II par un chargement
forfaitaire, peu sensible à la qualité réelle de l'infrastructure numérique
d'un assureur. Parallèlement, le règlement DORA impose depuis 2025 un cadre
détaillé de résilience opérationnelle numérique, sans pour autant fournir
de traduction directe en exigence de capital.

\paragraph{Objet du mémoire.} L'objet de ce mémoire est la
\textbf{quantification actuarielle du risque cyber} d'une entité financière,
envisagée non pas comme un exercice descriptif de conformité, mais comme un
exercice de mesure des \textbf{pertes potentielles prospectives}
susceptibles de résulter d'un non-respect, partiel ou total, des exigences du
règlement DORA. Il ne s'agit donc pas de vérifier, poste par poste, la
conformité réglementaire d'un assureur, mais de traduire un niveau de
non-conformité --- même hypothétique ou graduel --- en une distribution de
pertes financières futures, puis en un surcroît de capital de solvabilité
mesurable. Cette démarche prospective est ce qui distingue le présent travail
d'un audit de conformité classique : l'ambition est de répondre à la
question \emph{\og que pourrait-il en coûter, en fonds propres, de ne pas se
conformer à DORA ?\fg{}}, et non à la question \emph{\og l'entité est-elle
conforme aujourd'hui ?\fg{}}.

C'est dans cet interstice, entre un règlement de résilience opérationnelle
et un cadre de solvabilité financière, que se situe la problématique de ce
mémoire : \emph{comment transformer la non-conformité à DORA en une variable
quantitative susceptible d'alimenter une mesure interne de capital sous
Solvabilité~II ?} Il ne s'agit pas de proposer une nouvelle formule
réglementaire de Pilier~1, mais de construire, dans l'esprit du Pilier~2 et
de l'ORSA, un premier pont quantitatif entre les deux cadres. Cette
articulation n'a, à notre connaissance, pas encore été formalisée de façon
aussi directe dans la littérature actuarielle française : l'essentiel des
travaux traite soit du risque cyber comme risque assurantiel classique, soit
de DORA comme sujet de conformité, rarement des deux ensemble sous une même
mécanique de capital.

Le mémoire se déploie en cinq parties. La première pose le cadre
réglementaire et conceptuel dans lequel s'inscrit la démarche, de Bâle~III à
DORA, en passant par Solvabilité~II. La deuxième présente les données
mobilisées, leurs apports et leurs limites. La troisième détaille
l'architecture de modélisation retenue, fondée sur une approche par
distribution des pertes (LDA) enrichie de quatre briques. La quatrième
formalise les scénarios de conformité DORA et la construction du
$\Delta_{DORA}$, mesure directe du surcoût prospectif de la non-conformité.
La cinquième présente les résultats obtenus, leur interprétation actuarielle
et leurs implications pour le pilotage prudentiel.

\vspace{0.5em}
\noindent\rule{\textwidth}{0.8pt}
\vspace{0.8em}

% ==============================================================================
\section*{État de l'art : de l'assurabilité du risque cyber à la réglementation DORA}
\addcontentsline{toc}{section}{État de l'art}
% ==============================================================================

Avant de détailler la modélisation mathématique retenue pour quantifier le Capital de Solvabilité Requis (SCR) lié à DORA, il est indispensable de positionner ce travail dans la littérature actuarielle récente. Le risque cyber, de par sa nature évolutive et systémique, bouscule les paradigmes traditionnels de l'assurance non-vie et de la gestion des risques opérationnels. 

%--------------------------------------------------------------
\subsection{Les limites de l'assurabilité et la rareté des données}
%--------------------------------------------------------------

La littérature s'est d'abord concentrée sur la question fondamentale de l'assurabilité du risque cyber. Biener, Eling et Wirfs \citep{BienerElingWirfs2015} ont démontré que le risque cyber peine à satisfaire les critères classiques d'assurabilité, principalement en raison de l'asymétrie d'information et du risque d'accumulation extrême.

Face à ce constat, le principal obstacle à la quantification prudentielle reste le manque de bases de données de sinistralité exhaustives. Les travaux fondateurs de \textbf{Caroline Hillairet} et \textbf{Olivier Lopez} \citep{HillairetLopez2021} ont largement contribué à structurer l'exploitation des données publiques telles que la base \textit{Privacy Rights Clearinghouse} (PRC). Maillart et Sornette \citep{MaillartSornette2010} avaient déjà mis en évidence la nécessité d'utiliser la Théorie des Valeurs Extrêmes (EVT) pour traiter les montants de pertes, démontrant que le coût des méga-fuites de données suit des distributions à queues lourdes (notamment des lois de Pareto) dont l'espérance mathématique peut, dans certains régimes, diverger. Cette caractéristique justifie pleinement notre choix méthodologique d'appliquer une distribution de Pareto Généralisée (GPD) sur les bases PRC et OpRisk, tout en imposant des plafonds opérationnels pour stabiliser l'allocation de capital.

%--------------------------------------------------------------
\subsection{Modélisation de la contagion et dynamique systémique}
%--------------------------------------------------------------

Le deuxième grand axe de recherche actuarielle sur le cyber concerne la contagion. Contrairement aux risques climatiques (géographiquement circonscrits), le risque cyber s'affranchit des frontières physiques via les réseaux informatiques. 

Les travaux de \textbf{Boumezoued}, \textbf{Hillairet} et leurs co-auteurs ont exploré l'utilisation de modèles épidémiologiques et de processus de branchement pour quantifier cette contagion au sein des portefeuilles d'assurance, notamment le processus de Hawkes multivarié \citep{BessyRolandBoumezouedHillairet2021}, qui capture l'auto-excitation des cyberattaques (une faille \textit{zero-day} déclenchant une vague de sinistres rapprochés dans le temps), et le modèle de contagion à structure de réseau \citep{HillairetLopezDOultremontSpoorenberg2022}. Plus récemment, Hillairet, Lopez et Sopgoui \citep{HillairetLopezSopgoui2026} proposent un modèle SIR stochastique couplant contagion épidémique (processus de Cox) et croissance granulaire des firmes, appliqué aux attaques LockBit de 2024 sur un portefeuille de près de 3\,000 entreprises franciliennes ; ce cadre distingue explicitement contagion idiosyncratique et systémique, une distinction que nous retrouvons en~1.5.

Bien que les processus de Hawkes et les modèles SIR soient extrêmement performants pour la tarification dynamique et la modélisation fine de la propagation, ils s'avèrent très lourds à implémenter pour un calcul de SCR réglementaire (qui nécessite une vision statique à un an, la VaR 99,5~\%). C'est pourquoi, dans ce mémoire, nous avons fait le pont entre ces constats de contagion cyber et les modèles de risque de crédit. En nous inspirant de la dynamique de dépendance systémique de type Vasicek (Merton), nous modélisons la probabilité conditionnelle de défaut de conformité d'une entité face à un choc macro-cyber. Le remplacement de ce facteur unique par un processus de contagion de type SIR/Cox, plus fidèle à la dynamique réelle mais nécessitant une réévaluation dynamique du capital, constitue une perspective naturelle de raffinement (cf.~Conclusion).

%--------------------------------------------------------------
\subsection{Risques extrêmes et agrégation (Copules)}
%--------------------------------------------------------------

La question de l'agrégation des risques extrêmes est centrale dans le calcul du SCR. Les travaux de \textbf{Farkas}, Lopez et Thomas \citep{FarkasLopezThomas2021} sur le risque opérationnel soulignent la difficulté d'allouer du capital lorsque les différentes lignes de métier (ou briques de risque) sont soumises à des chocs lourds simultanés ; leur méthode d'arbres de régression Pareto généralisés, qui segmente automatiquement la sévérité par sous-population homogène plutôt que d'imposer un seuil global unique, rejoint l'esprit de la segmentation manuelle Systems Security / Business Disruption que nous opérons en~2.4.4 sur OpRisk Global --- sans toutefois en reproduire l'automatisation par arbre.

La formule standard de Solvabilité II utilise des matrices de corrélation linéaire qui échouent mathématiquement à capturer la dépendance dans les extrêmes (la queue de distribution). La littérature préconise l'usage de copules. Dans le sillage de ces recommandations, notre modèle rejette la copule Gaussienne au profit de la \textbf{copule de Gumbel}, garantissant une dépendance de queue supérieure stricte, essentielle pour simuler un scénario où une amende DORA majeure coïnciderait avec la défaillance systémique d'un prestataire cloud de rang 1.

%--------------------------------------------------------------
\subsection{De Solvabilité II à DORA : Le pont quantitatif}
%--------------------------------------------------------------

Historiquement, le risque cyber était relégué au sein du \og risque opérationnel\fg{}, quantifié de manière rudimentaire par la formule standard de Solvabilité II via un pourcentage des primes (BSCR). Cette approche \og forfaitaire\fg{} ne récompense pas les investissements en cybersécurité. 

Le règlement européen DORA (Digital Operational Resilience Act) change ce paradigme en imposant des standards stricts de résilience IT, de gestion des tiers prestataires et de report d'incidents. Cependant, DORA est un texte de conformité qualitative. Le défi actuel de la place financière est de traduire ces exigences en impacts bilanciels. 

C'est très exactement dans cette dynamique d'innovation méthodologique que s'inscrivent les travaux internes menés chez Nexialog Consulting sous la supervision de \textbf{Hugo Rapior}. L'approche développée dans ce mémoire formalise cette traduction : il s'agit de passer d'un constat binaire (conforme / non-conforme) à une modélisation continue du risque (le Delta DORA). En isolant la brique \og Prestataires\fg{} (Chapitre V de DORA) et la brique \og Sanctions\fg{} administratives, ce modèle interne partiel de type LDA propose un premier pont quantitatif rigoureux entre la résilience opérationnelle numérique (DORA) et la solvabilité prudentielle (Pilier 1 et ORSA de Solvabilité II).

%==============================================================
\section{Le cadre réglementaire : de Bâle III et Solvabilité II au règlement DORA}
%==============================================================

La quantification du capital requis au titre de la non-conformité au règlement DORA
s'inscrit dans un écosystème prudentiel dont il importe de retracer l'architecture.
Le risque que nous cherchons à modéliser --- le risque opérationnel d'origine
technologique, et le risque de manquement aux exigences de résilience numérique ---
se situe à l'intersection de trois cadres réglementaires construits successivement,
selon des philosophies distinctes mais convergentes : le dispositif bancaire de Bâle,
la directive Solvabilité~II propre à l'assurance, et le règlement transsectoriel DORA.
Cette section retrace cette construction, identifie le \emph{pont quantitatif} que DORA
permet d'établir entre exigence de capital et résilience opérationnelle, et formalise la
non-conformité comme une variable latente exploitable dans l'évaluation interne des risques.

%--------------------------------------------------------------
\subsection{L'écosystème réglementaire : Bâle III, Solvabilité II, DORA}
%--------------------------------------------------------------

\paragraph{Le risque opérationnel : une définition partagée, des traitements divergents.}
Le risque opérationnel a connu, au cours des deux dernières décennies, une mutation
profonde marquée par une recherche de convergence entre les secteurs bancaire et
assurantiel. Défini par le Comité de Bâle puis repris par la directive Solvabilité~II,
il désigne le risque de perte résultant de l'inadaptation ou de la défaillance de
procédures internes, de personnels et de systèmes, ou d'événements extérieurs. Cette
définition inclut le risque juridique mais exclut les risques stratégiques et de
réputation.

Si la définition est partagée, les modes de quantification ont divergé. Le secteur
bancaire a progressivement abandonné les modèles internes au profit d'une approche
standardisée, tandis que l'assurance a conservé une liberté méthodologique. Cette
asymétrie, que nous détaillons ci-après, éclaire le choix méthodologique du présent
mémoire : c'est précisément parce que Solvabilité~II laisse subsister un espace de
modélisation interne que la quantification du risque cyber sous forme d'une distribution
de pertes y conserve toute sa pertinence.

\paragraph{Bâle III/IV : du modèle interne AMA à la standardisation SMA.}
Sous le régime de Bâle~II, les établissements bancaires d'envergure internationale
étaient encouragés à développer des modèles internes complexes, regroupés sous
l'appellation d'approche de mesure avancée (\emph{Advanced Measurement Approaches},
AMA). Ces modèles reposaient principalement sur l'approche par distribution des pertes
(\emph{Loss Distribution Approach}, LDA) --- celle-là même que nous mobilisons dans ce
mémoire --- ou sur des approches par scénarios. S'ils offraient une grande flexibilité
théorique, ils ont engendré une variabilité excessive et difficilement justifiable des
actifs pondérés par le risque d'un établissement à l'autre, érodant la confiance des
superviseurs dans les ratios de solvabilité déclarés.

La finalisation des réformes de Bâle~III, communément désignée sous le nom de Bâle~IV,
a acté le retrait des modèles internes AMA pour le calcul du capital de Pilier~1, à
compter du 1\textsuperscript{er}~janvier~2022, au profit d'une méthode unique et
comparable : l'approche d'évaluation standardisée (\emph{Standardised Measurement
Approach}, SMA). L'exigence de capital pour risque opérationnel ($ORC$) s'exprime alors
comme le produit d'une composante de volume d'activité, le \emph{Business Indicator
Component} ($BIC$), et d'un facteur d'échelle fondé sur l'historique des pertes internes,
l'\emph{Internal Loss Multiplier} ($ILM$) :
\begin{equation}
ORC = BIC \times ILM,
\qquad
RWA_{Op} = 12{,}5 \times ORC.
\end{equation}

Le \emph{Business Indicator} synthétise la taille de l'établissement à partir de trois
composantes comptables (intérêts et dividendes, services, activité financière), auxquelles
sont appliqués des coefficients marginaux croissants par tranches. Le profil de risque
propre est réintroduit par l'$ILM$, qui compare le $BIC$ à la composante de pertes
historiques $LC$, égale à quinze fois la moyenne des pertes opérationnelles annuelles
significatives des dix dernières années. Lorsque la sinistralité historique est faible
relativement à l'activité ($LC < BIC$), l'$ILM$ est inférieur à l'unité et réduit
l'exigence ; lorsqu'elle est élevée, il agit comme une pénalité en fonds propres.

\paragraph{Enseignement pour notre démarche.} Le retrait de l'AMA n'invalide nullement
l'approche LDA en tant qu'outil de modélisation : il en restreint l'usage \emph{pour le
calcul réglementaire de Pilier~1 dans le secteur bancaire}. La logique de la composante
$LC$ --- faire dépendre le capital de l'historique réel des pertes --- demeure
parfaitement pertinente, et c'est cette logique, transposée à l'assurance et enrichie de
la théorie des valeurs extrêmes, que nous mettons en oeuvre.

\paragraph{Solvabilité II : l'asymétrie réglementaire et le maintien du choix méthodologique.}
Contrairement au secteur bancaire, le secteur de l'assurance conserve sous Solvabilité~II
une liberté de choix pour la quantification du risque opérationnel : formule standard,
paramètres propres à l'entreprise (USP), ou modèle interne, partiel ou global, approuvé
par l'autorité de contrôle. Dans la formule standard, l'exigence pour risque opérationnel
($SCR_{Op}$) s'ajoute de manière additive au capital de solvabilité requis de base
($BSCR$), ajusté de la capacité d'absorption des pertes :
\begin{equation}
SCR = BSCR + Adj + SCR_{Op},
\qquad
SCR_{Op} = \min(0{,}3 \times BSCR,\ Op) + 0{,}25 \times Exp_{ul},
\end{equation}
où $Op = \max(Op_{premiums}, Op_{provisions})$ est calculé à partir de coefficients
forfaitaires appliqués aux primes et provisions, vie et non-vie. Nous en donnons ici la
forme simplifiée ; la formule complète de l'Article~204 du règlement délégué~2015/35
comporte en outre des termes de croissance, sous la forme
$\max\!\left(0;\, 0{,}04 \times \Delta\!\right)$, pénalisant les hausses de primes
supérieures à $20\,\%$ d'un exercice à l'autre.

Cette approche présente une faiblesse conceptuelle majeure pour l'actuaire. En se fondant
exclusivement sur des indicateurs de volume, elle pose l'hypothèse implicite que deux
assureurs gérant le même volume de primes présentent une exposition identique au risque
opérationnel --- indépendamment de la qualité de leur infrastructure informatique, du
degré d'automatisation de leurs processus, ou de leur dépendance à des prestataires
d'externalisation critiques. Or c'est précisément cette dépendance technologique que DORA
vient encadrer, et dont nous cherchons à mesurer l'effet sur le capital.

\paragraph{Pilier 1 et Pilier 2.} Le $SCR_{Op}$ relève du Pilier~1, quantitatif et
forfaitaire. Mais Solvabilité~II comporte un second pilier, l'évaluation interne des
risques et de la solvabilité (\emph{Own Risk and Solvency Assessment}, ORSA), qui impose
à l'assureur une évaluation prospective de l'adéquation entre son profil de risque réel
et ses besoins de solvabilité. C'est dans le cadre de l'ORSA --- et non du forfait de
Pilier~1, plafonné à $30\,\%$ du $BSCR$ --- que notre modélisation trouve sa place
naturelle : l'ORSA autorise, et même encourage, le recours à des scénarios de stress
stochastiques que la formule standard ne sait pas produire.

\paragraph{DORA : un cadre transversal de résilience opérationnelle numérique.}
Face à l'insuffisance des exigences de capital de Pilier~1 pour garantir la continuité
des opérations en cas de crise informatique d'envergure, l'Union européenne a introduit
le règlement sur la résilience opérationnelle numérique du secteur financier
(\emph{Digital Operational Resilience Act}, règlement UE~2022/2554). Entré en application
le 17~janvier~2025, DORA agit comme une \emph{lex specialis} par rapport à la directive
NIS~2 et harmonise les obligations de gestion des risques liés aux technologies de
l'information et de la communication (TIC) à travers vingt-et-une catégories
d'entités financières : banques, entreprises d'assurance et de réassurance,
intermédiaires, ainsi que les prestataires tiers critiques de services TIC.

Le postulat de DORA est que la résilience opérationnelle ne saurait se réduire à
l'allocation de tampons de fonds propres. Le règlement impose un cadre proactif, articulé
autour de la capacité à \emph{résister, absorber, récupérer et s'adapter} face aux
perturbations informatiques majeures.

\paragraph{Position dans l'écosystème.} La progression Bâle~$\rightarrow$~Solvabilité~II
$\rightarrow$~DORA traduit un déplacement du centre de gravité prudentiel : de la
\emph{couverture financière} du risque (détenir assez de capital) vers la \emph{continuité
opérationnelle} (rester capable d'opérer). DORA ne remplace pas les exigences de capital ;
il en révèle l'angle mort, à savoir la dynamique propre des défaillances technologiques,
que ni le $SCR_{Op}$ forfaitaire ni l'ORSA qualitatif ne savaient capturer.

%--------------------------------------------------------------
\subsection{Ce que DORA exige concrètement : les cinq piliers ICT}
%--------------------------------------------------------------

Le règlement DORA structure ses exigences autour de cinq piliers, qui couvrent
l'ensemble du cycle de vie du risque informatique, de sa prévention à sa
résolution. Chacun d'entre eux introduit une granularité que ni la formule
standard, ni l'ORSA qualitatif traditionnel ne parvenaient à atteindre.

\paragraph{Pilier 1 --- Gouvernance et gestion du risque TIC.} DORA impose
l'existence d'un cadre de gouvernance formalisé, porté au plus haut niveau de
l'organe de direction, incluant une cartographie des actifs informationnels et
des dépendances critiques, une politique de sécurité documentée, et un
dispositif de contrôle interne capable d'identifier et de corriger les
faiblesses technologiques de façon continue --- non plus seulement lors
d'audits ponctuels.

\paragraph{Pilier 2 --- Gestion, classification et notification des incidents.}
Les entités doivent disposer d'un processus structuré de détection, de
qualification et de suivi des incidents liés aux TIC. Les incidents jugés
majeurs doivent être notifiés selon un calendrier resserré : une alerte sous
quatre heures après classification (et au plus tard vingt-quatre heures après
détection), un rapport intermédiaire sous soixante-douze heures, puis un
rapport final sous un mois.

\paragraph{Pilier 3 --- Tests de résilience opérationnelle numérique.} Il ne
suffit plus de documenter des procédures : elles doivent être éprouvées. Les
entités identifiées comme significatives sont soumises à des tests de
pénétration guidés par la menace (\emph{Threat-Led Penetration Testing},
TLPT), à mener au moins tous les trois ans, en conditions proches du réel.

\paragraph{Pilier 4 --- Gestion du risque lié aux prestataires tiers ICT.}
Ce pilier est central dans un environnement marqué par l'externalisation et le
recours au cloud. Il impose une cartographie exhaustive des contrats
critiques via un Registre d'Informations harmonisé, une gouvernance
contractuelle renforcée (clauses impératives de l'article~30), et une
surveillance directe des prestataires tiers critiques par les autorités
européennes, jusqu'au rang~3 de la chaîne de sous-traitance.

\paragraph{Pilier 5 --- Partage d'informations sur les cybermenaces.} DORA
encourage les entités à partager, de façon volontaire et encadrée, des
informations sur les menaces observées, dans une logique de résilience
collective. Ce pilier reste moins directement mobilisable dans le calibrage
actuariel du présent mémoire, mais il rappelle que DORA porte aussi une
dimension systémique de la résilience numérique, au-delà de la seule
entité individuelle.

\begin{center}
\renewcommand{\arraystretch}{1.3}
\begin{tabular}{>{\raggedright}p{3cm} >{\raggedright}p{5cm} >{\raggedright\arraybackslash}p{5.5cm}}
\toprule
\rowcolor{lightblue}
\textbf{Dimension} & \textbf{Solvabilité II (ORSA)} & \textbf{DORA} \\
\midrule
Taxonomie du risque & Traitement agrégé dans le module $SCR_{Op}$ &
Taxonomie granulaire des actifs TIC et indicateurs de risque distincts \\
Reporting des incidents & Notification interne, dialogue prudentiel non standardisé &
Notification obligatoire à l'ACPR : alerte sous 4\,h après classification comme majeur
(et au plus tard 24\,h après détection), rapport intermédiaire 72\,h,
rapport final 1~mois \\
Tests de résilience & Stress tests orientés adéquation du capital &
Tests de continuité et TLPT (pénétration guidée par la menace) au moins tous les 3~ans
pour les entités identifiées comme significatives \\
Cartographie des dépendances & Analyse générale des accords d'externalisation &
Registre d'Informations harmonisé, chaînes de sous-traitance jusqu'au rang~3 \\
Encadrement des tiers & Lignes directrices EIOPA non contraignantes &
Clauses contractuelles impératives (Art.~30), surveillance directe des prestataires
critiques \\
\bottomrule
\end{tabular}
\end{center}

Le \emph{Registre d'Informations} (ROI) mérite une attention particulière. En imposant
le recensement exhaustif et structuré des contrats TIC critiques et de leurs chaînes de
sous-traitance, DORA produit, pour la première fois, une donnée exploitable pour modéliser
quantitativement le risque de concentration technologique --- donnée qui faisait
précisément défaut à l'ORSA traditionnel.

%--------------------------------------------------------------
\subsection{1.3 --- ORSA $\leftrightarrow$ DORA : recouvrement et divergences}
%--------------------------------------------------------------

\paragraph{Les limites structurelles de l'ORSA traditionnel.} L'ORSA
traditionnel s'appuyait principalement sur des cartographies de risques
fondées sur l'avis d'experts. Trois limites en découlaient.

\textbf{Un cloisonnement méthodologique.} Les risques TIC étaient traités séparément des
risques financiers et assurantiels, sans intégration des corrélations réelles.
L'interaction entre une défaillance de l'outil de tarification (risque opérationnel) et
la sous-évaluation systématique des primes qui en résulte (risque de souscription) était
rarement modélisée de manière dynamique.

\textbf{L'absence de données de sinistralité qualifiées.} Les assureurs ne disposaient
pas de base standardisée permettant d'alimenter les modèles de sévérité pour les
événements cyber de faible fréquence et de forte gravité. L'étalonnage des distributions
à queue lourde restait hautement incertain et sujet à une forte volatilité de modèle ---
difficulté que notre travail aborde frontalement en mobilisant plusieurs sources
indépendantes (PRC, SAS~OpRisk~Global, Hackmageddon) et en assumant explicitement
l'incertitude résiduelle.

\textbf{La non-prise en compte des interdépendances de la chaîne de valeur.} Les
politiques d'externalisation sous Solvabilité~II se limitaient à des audits documentaires
ponctuels. L'effet de cascade d'une panne systémique chez un fournisseur d'infrastructure
cloud partagé par plusieurs acteurs du marché n'était pas intégré aux tests de résistance.

\paragraph{Zone de recouvrement.} ORSA et DORA partagent un même objet --- le risque opérationnel d'origine technologique
--- et une même exigence d'évaluation prospective. Trois domaines se recouvrent
explicitement : le risque opérationnel informatique, traité de façon agrégée par l'ORSA
et granulaire par DORA ; les scénarios de stress cyber, que l'ORSA peut désormais alimenter
des données et tests imposés par DORA ; et la continuité d'activité, exigence commune mais
de portée différente. C'est dans ce recouvrement que se loge la contribution du présent
mémoire : utiliser les artefacts produits par DORA (taxonomie, registre, tests) pour
\emph{quantifier} dans l'ORSA ce que celui-ci ne savait jusqu'alors que décrire.

\paragraph{Zones de divergence fondamentale.}
\textbf{La finalité de l'évaluation.} Solvabilité~II s'inscrit dans une logique de
continuité \emph{financière} : prouver que l'assureur détient assez de fonds propres pour
absorber les pertes d'une défaillance majeure. DORA s'inscrit dans une logique de
continuité \emph{opérationnelle} : garantir que, même sans capital excédentaire,
l'infrastructure maintient les services critiques et se restaure rapidement. Les deux
cadres sont complémentaires, non substituables.

\textbf{Le principe de proportionnalité.} Solvabilité~II applique une proportionnalité
fondée sur la taille globale (primes, provisions). DORA décline la sienne selon la
complexité des systèmes d'information, la dépendance au cloud tiers, et la criticité dans
l'écosystème de distribution. Une captive de réassurance ou une petite mutuelle peut ainsi
bénéficier d'un régime DORA simplifié sans pour autant échapper aux exigences de capital.

\textbf{La temporalité de la supervision.} Solvabilité~II fonctionne sur des cycles
trimestriels et annuels (QRT, SFCR, RSR). DORA impose une réactivité quasi temps réel,
notamment par la notification des incidents majeurs sous quatre heures, contraignant
l'assureur à des processus permanents de veille connectés aux canaux du régulateur.

%--------------------------------------------------------------
\subsection{1.4 --- La probabilité de défaut de conformité comme variable latente}
%--------------------------------------------------------------

Les sections précédentes ont établi que DORA fournit des artefacts nouveaux ---
registre des tiers, notifications d'incidents, résultats de tests TLPT ---
mais aucun de ces artefacts ne constitue, pris isolément, une mesure directe
du niveau de conformité d'une entité. Une entité peut être en avance sur la
gouvernance TIC mais en retard sur les tests de résilience ; une autre peut
être conforme sur le papier mais fortement exposée via un prestataire tiers
critique mal maîtrisé. La conformité DORA n'est donc pas une variable
observable au sens strict : c'est une caractéristique multidimensionnelle,
partiellement cachée, qu'il convient de formaliser par une \textbf{variable
latente à structure factorielle}, sur le modèle du cadre de Merton-Vasicek
utilisé pour la modélisation des défauts corrélés en risque de crédit --- et
qui constitue également le fondement de l'approche IRB de Bâle pour la
corrélation des probabilités de défaut.

\paragraph{Structure à facteur unique.} Soit $Z_i$ la variable latente
continue de fragilité réglementaire et opérationnelle de l'entité $i$,
agrégeant de façon non observable ses faiblesses sur les cinq piliers DORA.
Suivant la structure de Vasicek, $Z_i$ se décompose en un facteur
systémique $\theta$, commun à l'ensemble du marché, et un facteur
idiosyncratique $\varepsilon_i$, propre à l'entité :
\begin{equation}
Z_i = \sqrt{\rho_i}\,\theta + \sqrt{1-\rho_i}\,\varepsilon_i,
\qquad \theta \sim \mathcal{N}(0,1), \quad \varepsilon_i \sim \mathcal{N}(0,1),
\end{equation}
avec $\theta$ et $\varepsilon_i$ indépendants, et $\rho_i \in [0,1]$ le
coefficient de corrélation de l'entité $i$ à l'environnement cyber
systémique. Un $\rho_i$ élevé caractérise une entité dont la posture de
conformité se dégrade fortement lorsque l'environnement sectoriel se
détériore (forte dépendance aux tiers, mutualisation d'infrastructures
critiques), tandis qu'un $\rho_i$ faible caractérise une entité dont la
gouvernance reste robuste indépendamment du contexte général.

\paragraph{Seuil de défaut de conformité.} Comme dans le modèle de Merton, le
défaut de conformité de l'entité $i$ est déclenché lorsque sa variable
latente franchit un seuil $K_i$, calibré de sorte à retrouver le taux de
non-conformité moyen observé sur le marché. La probabilité de défaut de
conformité conditionnelle à l'état de l'environnement systémique $\theta$
s'écrit alors :
\begin{equation}
PCD_i(\theta) = \mathbb{P}\bigl(Z_i \leq K_i \,\big|\, \theta\bigr)
= \Phi\!\left(\frac{K_i - \sqrt{\rho_i}\,\theta}{\sqrt{1-\rho_i}}\right),
\end{equation}
où $\Phi$ désigne la fonction de répartition de la loi normale centrée
réduite. Le seuil $K_i$ est obtenu par inversion de la probabilité de défaut
inconditionnelle $\bar{PCD}_i$ (calibrée sur les taux de non-conformité
observés en 2025) : $K_i = \Phi^{-1}(\bar{PCD}_i)$.

Cette écriture retrouve les propriétés attendues : lorsque l'environnement
systémique se dégrade ($\theta \to -\infty$, interprété comme un choc cyber
sectoriel négatif), $PCD_i(\theta) \to 1$ pour toute entité présentant
$\rho_i > 0$ ; à l'inverse, une entité affichant $\rho_i \to 0$ conserve un
$PCD_i(\theta)$ proche de son niveau inconditionnel $\bar{PCD}_i$, quasiment
insensible au contexte de marché. Le paramètre $\rho_i$ joue ainsi le rôle
d'un coefficient de dépendance à la contagion systémique, au sens strict du
modèle à facteur unique de Vasicek.

\paragraph{Interpolation entre scénarios-types.} Cette formalisation permet de
dépasser le choix discret d'un scénario de marché $S_x \in \{S_0, S_1, S_2\}$
(conforme, partiellement conforme, non conforme) au profit d'une lecture
continue : $S_0$ correspond au cas limite $PCD_i(\theta) \to 0$, $S_2$ au cas
limite $PCD_i(\theta) \to 1$, et $S_1$ à une valeur intermédiaire calibrée sur
les profils de conformité partielle observés en 2025. Le multiplicateur de
fréquence appliqué à un vecteur d'attaque $v$ pour l'entité $i$ s'écrit alors
comme une interpolation linéaire entre le régime conforme et le régime de
non-conformité totale sourcé empiriquement :
\begin{equation}
m_v\bigl(PCD_i(\theta)\bigr) = 1 + PCD_i(\theta) \times \left(\bar{m}_v^{S_2} - 1\right),
\end{equation}
où $\bar{m}_v^{S_2}$ est le centre de l'intervalle de multiplicateur sourcé en
scénario de non-conformité totale pour le vecteur $v$ (cf. Partie~4). Cette
écriture ne requiert aucune calibration supplémentaire au-delà de $\rho_i$ et
$\bar{PCD}_i$ : elle réutilise directement les multiplicateurs déjà sourcés
par vecteur d'attaque, en substituant au choix manuel d'un scénario de marché
la probabilité propre à chaque profil d'entité, conditionnellement à l'état
de l'environnement cyber systémique $\theta$.

Trois avantages méthodologiques découlent de ce choix de structure. Premièrement,
le cadre de Vasicek est un standard éprouvé de la modélisation actuarielle et
prudentielle des risques corrélés, ce qui confère au $PCD_i(\theta)$ une
assise théorique plus solide qu'une fonction de lien ad hoc. Deuxièmement, le
paramètre $\rho_i$ offre une lecture économique directe et interprétable de
la sensibilité de chaque entité à la contagion systémique --- une lecture
qui prépare naturellement la discussion de la contagion interne et externe
menée en 1.5. Troisièmement, la structure à facteur unique est immédiatement
compatible avec l'architecture de dépendance retenue en Partie~3 (copule de
Gumbel et modèle à facteur commun), puisqu'elle repose sur le même principe
de décomposition entre un choc commun $\theta$ et des chocs idiosyncratiques
$\varepsilon_i$.

Cette variable latente n'est pas un simple artifice de modélisation : elle
constitue le point de jonction technique entre le cadre réglementaire décrit
ci-dessus et le moteur actuariel développé dans la suite du mémoire. C'est
elle qui, dans le chapitre consacré à la fréquence, vient moduler l'intensité
des sinistres simulés et piloter la probabilité d'occurrence d'une sanction
réglementaire.

%--------------------------------------------------------------
\subsection{1.5 --- Contagion interne et contagion externe via tiers ICT}
%--------------------------------------------------------------

L'introduction d'une variable latente de conformité invite à distinguer deux
canaux de propagation du risque, que le cadre DORA rend précisément
observables : une contagion \textbf{interne} et une contagion
\textbf{externe}.

\paragraph{Contagion interne.} Une faiblesse de gouvernance TIC, un retard
dans le déploiement des correctifs de sécurité, ou une insuffisance de
sensibilisation du personnel ne restent pas cantonnés à un seul type
d'incident. Ils dégradent simultanément plusieurs vecteurs d'attaque : un
défaut de formation accroît la vulnérabilité au phishing, mais fragilise
également la détection précoce d'autres intrusions. Cette dimension justifie
que la fréquence des incidents ne soit pas traitée comme la simple somme de
processus indépendants, mais comme le produit d'un facteur commun de
maîtrise opérationnelle interne.

\paragraph{Contagion externe.} Le pilier consacré à la gestion du risque lié
aux prestataires tiers ICT rend cette seconde dimension explicite. Une entité
peut afficher une gouvernance interne robuste tout en restant exposée à la
défaillance d'un prestataire critique --- hébergeur cloud, éditeur logiciel,
sous-traitant de paiement --- partagé par plusieurs acteurs du marché. Le
Registre d'Informations et la traçabilité des chaînes de sous-traitance
jusqu'au rang~3, imposés par DORA, rendent pour la première fois cette
dépendance cartographiable, alors qu'elle échappait largement aux audits
documentaires ponctuels de Solvabilité~II.

Cette distinction entre contagion interne et externe n'est pas seulement
conceptuelle : elle structure directement l'architecture de modélisation
retenue dans ce mémoire. La contagion interne motive le recours à une
structure de dépendance de type facteur commun entre les briques de perte
liées à la maîtrise opérationnelle de l'entité elle-même. La contagion
externe justifie l'introduction d'une brique de risque dédiée au prestataire,
distincte de la brique de remédiation, et calibrée sur un scénario
d'occurrence propre --- car le risque tiers ne se comporte pas comme un
sinistre interne ordinaire, mais comme un choc importé, potentiellement
corrélé entre plusieurs entités du marché exposées au même fournisseur.

%--------------------------------------------------------------
\subsection*{Positionnement du mémoire}
%--------------------------------------------------------------

Au terme de ce panorama, la contribution visée se précise. Le forfait de Pilier~1
($SCR_{Op}$, plafonné à $30\,\%$ du $BSCR$) est, par construction, incapable de refléter
le risque de cyber-catastrophe systémique : il ignore la qualité de l'infrastructure, la
dépendance aux tiers, et la dynamique propre des défaillances. Nous proposons donc, dans
la logique du Pilier~2, un modèle de stress stochastique fondé sur la théorie des valeurs
extrêmes, étalonné sur des données de sinistralité sectorielles et structuré par la
variable latente de conformité introduite ci-dessus. L'objectif n'est pas de produire un
chiffre unique de capital, mais une \emph{distribution} de l'exigence, assortie de ses
intervalles d'incertitude --- car le SCR cyber n'est pas un nombre, mais une distribution
large.

% ==============================================================================
\section*{Cadre mathématique de la modélisation des risques extrêmes}
\addcontentsline{toc}{section}{Cadre mathématique de la modélisation des risques extrêmes}
% ==============================================================================

La quantification du risque cyber, caractérisée par une asymétrie prononcée et la présence de sinistres d'une sévérité exceptionnelle, nécessite un cadre probabiliste robuste. Les lois de probabilité traditionnelles (Lognormale, Gamma, Weibull), souvent utilisées en assurance non-vie pour modéliser le corps de la distribution (la fréquence des sinistres d'attrition), échouent à capturer l'épaisseur des queues de distribution typiques des méga-fuites de données ou des paralysies systémiques. Cette section détaille les fondements théoriques de la Théorie des Valeurs Extrêmes (EVT) sous-jacents au modèle de distribution de pertes agrégées (LDA) implémenté, ainsi que la modélisation de la dépendance.

%--------------------------------------------------------------
\subsection{Justification de l'EVT et Théorème de Fisher-Tippett-Gnedenko (GEV)}
%--------------------------------------------------------------

Là où le Théorème Central Limite s'intéresse au comportement asymptotique des sommes de variables aléatoires, la Théorie des Valeurs Extrêmes (EVT) étudie le comportement asymptotique de leurs maxima. 

Soit une séquence $(X_1, X_2, \dots, X_n)$ de variables aléatoires indépendantes et identiquement distribuées (i.i.d.) de fonction de répartition $F$, représentant la sévérité des cyberattaques. L'approche initiale de l'EVT, dite des \textit{Block Maxima}, s'intéresse au maximum sur des blocs de taille $n$ (par exemple, le sinistre cyber maximal annuel) :
\begin{equation}
M_n = \max(X_1, X_2, \dots, X_n)
\end{equation}

\paragraph{Théorème de Fisher-Tippett-Gnedenko (1928, 1943).} S'il existe des suites de constantes de normalisation $a_n > 0$ et $b_n \in \mathbb{R}$, et une distribution non dégénérée $H$ telles que :
\begin{equation}
\lim_{n \to \infty} \mathbb{P}\left( \frac{M_n - b_n}{a_n} \le x \right) = \lim_{n \to \infty} F^n(a_n x + b_n) = H(x)
\end{equation}
Alors, $H$ appartient nécessairement à la famille des distributions des extrêmes généralisées (GEV - \textit{Generalized Extreme Value}), dont la fonction de répartition s'écrit :
\begin{equation}
H_{\xi}(x) = 
\begin{cases} 
\exp\left(-\left(1 + \xi x \right)^{-1/\xi}\right) & \text{si } \xi \neq 0 \\
\exp\left(-\exp(-x)\right) & \text{si } \xi = 0 
\end{cases}
\end{equation}
définie sur l'ensemble $\{x : 1 + \xi x > 0\}$. Le paramètre $\xi$ (indice de queue) définit trois grands domaines d'attraction :
\begin{itemize}
    \item \textbf{$\xi > 0$ (Loi de Fréchet) :} La distribution présente une queue lourde à décroissance polynomiale. C'est le régime critique qui nous intéresse pour les sinistres cyber majeurs.
    \item \textbf{$\xi = 0$ (Loi de Gumbel) :} La distribution présente une queue légère à décroissance exponentielle (domaine d'attraction des lois Normale, Lognormale, Gamma).
    \item \textbf{$\xi < 0$ (Loi de Weibull) :} La distribution a un support fini supérieur (queue courte).
\end{itemize}

%--------------------------------------------------------------
\subsection{Le passage au Peaks-Over-Threshold (POT) et la GPD}
%--------------------------------------------------------------

Bien que rigoureuse, l'approche \textit{Block Maxima} (GEV) présente une limite majeure en actuariat pratique : elle induit un immense gaspillage d'information. En ne retenant que le maximum de chaque bloc, elle occulte potentiellement le deuxième plus grand sinistre d'une année extrême, qui pourrait pourtant être supérieur au maximum d'une année calme. Dans le contexte du risque cyber, où les données historiques fiables et quantifiées sont particulièrement rares, un tel gaspillage statistique fragiliserait la modélisation.

Pour contourner ce problème, ce mémoire privilégie l'approche \textit{Peaks-Over-Threshold} (POT). Au lieu de ne garder que les maxima, cette méthode exploite \textbf{toutes} les observations excédant un seuil élevé $u$. On s'intéresse alors à la distribution des excès $Y = X - u$.

La fonction de répartition conditionnelle des excès est donnée par :
\begin{equation}
F_u(y) = \mathbb{P}(X - u \le y \mid X > u) = \frac{F(u+y) - F(u)}{1 - F(u)}
\end{equation}

\paragraph{Théorème de Balkema-de Haan-Pickands (1974).} Ce théorème établit le pont mathématique entre la GEV et l'approche POT. Pour une large classe de distributions sous-jacentes $F$ (celles appartenant au domaine d'attraction d'une loi GEV $H_\xi$), il existe une fonction mesurable positive $\sigma(u)$ telle que, lorsque le seuil $u$ tend vers la limite supérieure du support de $F$ ($x_F$), la distribution des excès converge asymptotiquement vers une distribution de Pareto Généralisée (GPD) :
\begin{equation}
\lim_{u \to x_F} \sup_{0 \le y < x_F - u} \left| F_u(y) - G_{\xi, \sigma(u)}(y) \right| = 0
\end{equation}

Où la fonction de répartition de la GPD s'écrit :
\begin{equation}
G_{\xi, \sigma}(y) = 
\begin{cases} 
1 - \left(1 + \frac{\xi y}{\sigma}\right)^{-1/\xi} & \text{si } \xi \neq 0 \\
1 - \exp\left(-\frac{y}{\sigma}\right) & \text{si } \xi = 0 
\end{cases}
\end{equation}
avec $y \ge 0$ si $\xi \ge 0$, et $0 \le y \le -\sigma/\xi$ si $\xi < 0$. Le paramètre de forme $\xi$ de la GPD est strictement le même que celui de la loi GEV correspondante. Dans le contexte du risque cyber, un indice de queue $\xi > 0$ traduit une distribution à queue lourde de type Fréchet, impliquant des moments potentiellement infinis (l'espérance est infinie si $\xi \ge 1$, la variance si $\xi \ge 0.5$).

%--------------------------------------------------------------
\subsection{Estimation des paramètres : Maximum de Vraisemblance vs Estimateur de Hill}
%--------------------------------------------------------------

L'estimation de l'indice de queue $\xi$ est un enjeu critique, car une infime variation de ce paramètre modifie massivement le Capital de Solvabilité Requis (SCR). Deux approches principales s'affrontent dans la littérature.

\paragraph{L'estimateur de Hill.}
Pour $\xi > 0$, l'estimateur non paramétrique de Hill (1975) repose sur les statistiques d'ordre. Soit $X_{(1)} \ge X_{(2)} \ge \dots \ge X_{(n)}$ l'échantillon ordonné de manière décroissante. L'estimateur de Hill fondé sur les $k$ plus grandes observations est défini par :
\begin{equation}
\hat{\xi}_{Hill}(k) = \frac{1}{k} \sum_{i=1}^{k} \ln X_{(i)} - \ln X_{(k+1)}
\end{equation}
Bien que cet estimateur soit asymptotiquement convergent et sans biais sous la stricte condition que la distribution sous-jacente appartienne exactement au domaine d'attraction de Fréchet, il présente en pratique un biais d'échantillon fini majeur. En effet, si la composante à variation lente de la queue de distribution ne converge pas assez vite, $\hat{\xi}_{Hill}$ devient instable.

\paragraph{Le Maximum de Vraisemblance (MLE).}
Pour pallier le biais de l'estimateur de Hill hors du régime asymptotique strict, ce mémoire privilégie l'estimation par Maximum de Vraisemblance (MLE) sur les excès de la GPD. La log-vraisemblance pour un échantillon de $N_u$ excès $y_i$ au-delà du seuil $u$ s'écrit (pour $\xi \neq 0$) :
\begin{equation}
\ell(\xi, \sigma) = -N_u \ln(\sigma) - \left(1 + \frac{1}{\xi}\right) \sum_{i=1}^{N_u} \ln\left(1 + \frac{\xi y_i}{\sigma}\right)
\end{equation}
La maximisation numérique de cette fonction fournit des estimateurs $(\hat{\xi}_{MLE}, \hat{\sigma}_{MLE})$ qui, bien que nécessitant des méthodes d'optimisation (telles que l'algorithme de Nelder-Mead utilisé dans nos scripts Python), offrent une robustesse supérieure face à la variance d'échantillonnage des bases de sinistres cyber comme OpRisk ou PRC.

%--------------------------------------------------------------
\subsection{Modélisation de la dépendance de queue : le choix de la Copule de Gumbel}
%--------------------------------------------------------------

L'agrégation des différentes briques de risque (remédiation, tiers prestataires, amendes) nécessite de capturer la structure de dépendance. Le théorème de Sklar (1959) permet de dissocier les lois marginales de la structure de dépendance.

\paragraph{Théorème de Sklar.} Soit $H$ une fonction de répartition jointe de marges $F_1, \dots, F_d$. Il existe une copule $C : [0,1]^d \to [0,1]$ telle que pour tout $x_1, \dots, x_d \in \mathbb{R} \cup \{-\infty, +\infty\}$ :
\begin{equation}
H(x_1, \dots, x_d) = C(F_1(x_1), \dots, F_d(x_d))
\end{equation}

Dans le cas du risque cyber, la préoccupation majeure du régulateur (notamment sous DORA) est la contagion systémique : la probabilité que les pires sinistres (ex: paralysie d'un prestataire cloud critique) surviennent simultanément avec des attaques majeures sur l'assureur. Cette caractéristique impose l'utilisation d'une copule présentant une \textbf{dépendance de queue supérieure} (upper tail dependence) strictement positive.

Le coefficient de dépendance de queue supérieure $\lambda_U$ entre deux variables $X_1, X_2$ de fonctions de répartition $F_1, F_2$ est défini par :
\begin{equation}
\lambda_U = \lim_{q \to 1^-} \mathbb{P}(X_1 > F_1^{-1}(q) \mid X_2 > F_2^{-1}(q))
\end{equation}
La copule Gaussienne possède une dépendance de queue asymptotiquement nulle ($\lambda_U = 0$), ce qui la disqualifie pour modéliser le risque de crise systémique. C'est pourquoi nous avons opté pour la famille archimédienne, et spécifiquement la \emph{copule de Gumbel}, définie pour $\eta \in [1, +\infty)$ par :
\begin{equation}
C_\eta(u, v) = \exp \left( - \left[ (-\ln u)^\eta + (-\ln v)^\eta \right]^{1/\eta} \right)
\end{equation}
Pour la copule de Gumbel, $\lambda_U = 2 - 2^{1/\eta}$. Le paramètre $\eta$ permet ainsi de piloter finement la propension des briques de risque à générer des pertes extrêmes conjointes.

%--------------------------------------------------------------
\subsection{Transformation d'un choc systémique en risque de conformité : le modèle de Vasicek}
%--------------------------------------------------------------

Pour modéliser l'occurrence de l'amende administrative et la défaillance d'un prestataire sous DORA en fonction d'un environnement macro-cyber (crise systémique), nous adaptons le modèle à un facteur de Vasicek (initialement conçu pour le risque de crédit asymétrique de type Merton).

Soit une variable latente $Z_i$ représentant la robustesse globale (ou solvabilité numérique) d'une entité $i$. Nous postulons que cette robustesse dépend d'un facteur de risque systémique commun à tout le marché $Y$ (par exemple, la découverte d'une faille \textit{zero-day} mondiale) et d'un facteur idiosyncratique propre à l'entité $X_i$ (la qualité de sa gouvernance IT) :
\begin{equation}
Z_i = \sqrt{\rho}Y + \sqrt{1-\rho}X_i
\end{equation}
où $Y$ et $X_i$ sont des variables normales centrées réduites indépendantes, et $\rho \in [0,1]$ mesure la sensibilité de l'entité au choc systémique (le coefficient de corrélation des actifs). 

Une entité subit une défaillance de conformité (déclenchant potentiellement une amende ou une rupture de service critique) si sa robustesse $Z_i$ tombe sous un certain seuil de tolérance $c_i = \Phi^{-1}(PCD_{marginal})$. Conditionnellement à un état donné de l'environnement macro-cyber $Y = y$ (où $y < 0$ représente un choc dégradé), la probabilité conditionnelle de défaut de conformité, notée $PCD_i(y)$, s'écrit :
\begin{equation}
PCD_i(y) = \mathbb{P}(Z_i < c_i \mid Y = y) = \mathbb{P}\left( \sqrt{\rho}y + \sqrt{1-\rho}X_i < c_i \right)
\end{equation}
Puisque $X_i \sim \mathcal{N}(0,1)$, on isole $X_i$ pour obtenir la formule fermée :
\begin{equation}
PCD_i(y) = \Phi \left( \frac{\Phi^{-1}(PCD_{marginal}) - \sqrt{\rho}y}{\sqrt{1-\rho}} \right)
\end{equation}
Cette formulation mathématique rigoureuse justifie la dynamique asymétrique observée dans nos résultats : une entité \og Leader\fg{} (dont le seuil $c_i$ est très bas) restera résiliente même si $Y$ se dégrade, tandis qu'une entité \og Médiane\fg{} verra sa probabilité de défaut converger asymptotiquement vers 1, ce qui active mécaniquement la contagion entre la brique prestataire et la brique sanction dans l'allocation d'Euler.

%==============================================================
\section{Données et limites}
%==============================================================

La modélisation développée dans ce mémoire ne repose sur aucune base de
données unique. Chaque source mobilisée présente un périmètre, un biais et un
niveau de granularité distincts, et c'est précisément leur articulation ---
plutôt que le choix d'une source \og privilégiée\fg{} --- qui constitue la
démarche méthodologique retenue. Cette partie présente successivement les
trois sources utilisées, avant de formaliser un cadre épistémique explicite
sur les limites de l'exercice.

%--------------------------------------------------------------
\subsection*{Bloc synthèse --- Cartographie des sources et de leurs rôles}
%--------------------------------------------------------------

Un point de clarification méthodologique central structure l'ensemble de
cette partie : \textbf{aucune base n'est utilisée seule}. Chaque source
alimente une composante précise du modèle, selon ses propriétés propres.

\begin{center}
\renewcommand{\arraystretch}{1.3}
\begin{tabular}{>{\raggedright}p{3cm} >{\raggedright}p{5.5cm} >{\raggedright\arraybackslash}p{5.5cm}}
\toprule
\rowcolor{lightblue}
\textbf{Source} & \textbf{Rôle dans le modèle} & \textbf{Justification / limite} \\
\midrule
PRC 2025 & Fréquence ($\lambda$) ; sévérité indirecte via conversion Jacobs &
Périmètre défini, mais sévérité dérivée d'un modèle log-log \\
\addlinespace
SAS OpRisk Global & Sévérité (montants réels) &
Montants observés, mais biais vers les grandes entités financières \\
\addlinespace
Hackmageddon & Multiplicateurs de scénario DORA (structure du risque) &
Ne calibre pas le modèle : ajuste $\lambda$ selon le niveau de conformité \\
\bottomrule
\end{tabular}
\end{center}

Cette répartition des rôles évite un écueil fréquent dans la modélisation du
risque cyber : traiter une base généraliste comme si elle pouvait, à elle
seule, calibrer l'ensemble d'un modèle fréquence-sévérité. Chaque source est
au contraire mobilisée là où elle est la plus fiable, et ses limites propres
sont explicitement documentées dans les sous-sections suivantes.

%--------------------------------------------------------------
\subsection{Exigences de Solvabilité II sur la qualité des données (Article 19)}
%--------------------------------------------------------------

L'utilisation d'un modèle interne, même partiel, est strictement encadrée par la directive Solvabilité II. L'Article 19 du Règlement Délégué (UE) 2015/35 impose que les données utilisées pour le calcul de la distribution de probabilité prévisionnelle respectent trois critères cumulatifs : l'exactitude, l'exhaustivité et la pertinence. 

Dans le domaine du risque cyber, aucune base de données publique ou commerciale ne satisfait pleinement ces trois critères. L'actuaire est donc contraint de documenter rigoureusement les biais inhérents aux données historiques et de justifier les retraitements appliqués. La validité du SCR calculé ne repose pas sur la perfection de la donnée (qui n'existe pas), mais sur la transparence des limites épistémiques et la robustesse du modèle mathématique face à ces biais.

%--------------------------------------------------------------
\subsection{Analyse critique des bases de sinistralité : les trois biais fondamentaux}
%--------------------------------------------------------------

L'utilisation conjointe des bases \textit{Privacy Rights Clearinghouse} (PRC) et SAS OpRisk Global permet de croiser les perspectives, mais nécessite d'identifier trois biais statistiques majeurs qui altèrent l'estimation de l'indice de queue $\xi$.

\paragraph{1. Le biais de sélection (Selection Bias).}
Ce biais affecte particulièrement la base OpRisk Global, qui se nourrit principalement d'événements rendus publics dans la presse financière ou les rapports annuels. Par construction, un incident cyber n'atteint ce seuil de visibilité que s'il frappe une entité d'une certaine taille (généralement des institutions systémiques ou de grandes banques internationales) ou si son montant est exceptionnellement élevé. 
Ce biais de taille (size bias) conduit à une sur-représentation des événements extrêmes et une sous-représentation de la sinistralité d'attrition. Si l'on calibre directement un modèle sur cette base pour une entité de taille moyenne (ex: une mutuelle régionale), la probabilité de subir une perte extrême sera mécaniquement surestimée. C'est pour cette raison que notre approche conditionne la fréquence d'occurrence au profil de l'entité via le modèle de Vasicek (Partie 3).

\paragraph{2. Le biais de déclaration (Reporting Bias).}
La base PRC est soumise à un biais de déclaration lié aux asymétries réglementaires. Aux États-Unis, la notification des fuites de données personnelles (Data Breaches) est obligatoire dans la plupart des États, ce qui rend la base PRC très riche sur ce vecteur d'attaque. En revanche, un rançongiciel (Ransomware) qui paralyse les serveurs de l'entreprise mais n'exfiltre aucune donnée personnelle peut ne déclencher aucune obligation de notification légale (avant l'entrée en vigueur d'actes récents comme le CIRCIA). 
Par conséquent, calibrer une fréquence globale sur la PRC tend à sous-estimer le risque d'interruption d'activité (Business Interruption) au profit du seul risque de compromission de données. C'est précisément pour corriger ce biais que nous avons superposé la structure des vecteurs d'attaque de la base Hackmageddon (qui capte l'ensemble du spectre, y compris les ransomwares) sur la fréquence de base issue de la PRC.

\paragraph{3. Le biais de survie (Survivorship Bias).}
Dans les bases de pertes opérationnelles, la sévérité maximale observée est paradoxalement bornée par la capacité de survie de l'entité. Une cyberattaque ayant entraîné la faillite pure et simple de l'entreprise (ruine) disparaît souvent des bases de pertes opérationnelles classiques, car il n'y a plus d'entité survivante pour reporter la perte financière finale et les coûts de remédiation à long terme. Ce biais tronque artificiellement l'extrême queue de distribution. La Théorie des Valeurs Extrêmes (EVT), en ajustant une loi de Pareto Généralisée (GPD) qui permet d'extrapoler au-delà du maximum historique observé, est la réponse mathématique adéquate à ce biais de survie.

%--------------------------------------------------------------
\subsection{Traitement de l'inflation cyber et actualisation des pertes}
%--------------------------------------------------------------

L'hétérogénéité temporelle des bases de données constitue une autre infraction potentielle au critère d'exactitude de Solvabilité II. La base OpRisk Global agrège des pertes allant de l'année 2000 à 2026. Utiliser un montant de perte de 2018 tel quel dans un modèle de capital en 2026 constitue une erreur actuarielle sévère, l'inflation cyber ne se comportant pas comme l'inflation macroéconomique classique (IPC).

Le coût d'un sinistre cyber est principalement tiré par trois facteurs :
\begin{itemize}
    \item \textbf{L'inflation salariale IT :} Le coût des experts en remédiation (forensic, restauration de systèmes) connaît une inflation nettement supérieure à l'indice des prix à la consommation, en raison de la pénurie de talents en cybersécurité.
    \item \textbf{L'inflation juridique et réglementaire :} L'entrée en vigueur du RGPD en 2018 a drastiquement modifié l'échelle des sanctions. Une fuite de données en 2015 coûtait infiniment moins cher en amendes administratives que la même fuite en 2025.
    \item \textbf{L'escalade des rançons :} La professionnalisation du cybercrime (Ransomware-as-a-Service) a transformé les demandes de rançons de quelques milliers de dollars en 2016 à plusieurs dizaines de millions d'euros en 2026.
\end{itemize}

Pour respecter les exigences de l'Article 19, les sinistres historiques doivent être actualisés en euros de l'année de modélisation (2026). Soit $L_t$ la perte observée l'année $t$, la perte actualisée $\tilde{L}_{2026}$ s'écrit :
\begin{equation}
\tilde{L}_{2026} = L_t \times \prod_{k=t}^{2025} (1 + i_k^{cyber})
\end{equation}
où $i_k^{cyber}$ est un taux d'inflation spécifique aux sinistres technologiques (souvent estimé par des indices sectoriels entre 5\,\% et 10\,\% annuels). Dans le cadre de ce mémoire, nous assumons que les bases commerciales traitées intègrent une homogénéisation monétaire, mais soulignons que la construction d'un indice d'inflation cyber robuste demeure un défi majeur pour la place financière.

Par ailleurs, le traitement des données manquantes dans la PRC (qui recense parfois le nombre d'enregistrements compromis sans le montant financier) est géré par la conversion de Jacobs. Bien qu'utile comme \textit{proxy}, ce modèle de régression log-log présente la limite de \og lisser\fg{} la variance et de minorer les extrêmes, justifiant notre décision de ne pas plafonner la calibration issue directement d'OpRisk (où les montants sont observés de bout en bout).

%--------------------------------------------------------------
\subsection{Perspectives : Le registre européen DORA comme solution pérenne}
%--------------------------------------------------------------

Les limites épistémiques et statistiques exposées ci-dessus contraignent aujourd'hui le calcul du SCR cyber à de larges intervalles d'incertitude. Cependant, cette situation est transitoire. L'apport fondamental du règlement DORA ne réside pas uniquement dans l'imposition de normes de sécurité (Chapitre II), mais dans la création d'un cadre de \textbf{notification harmonisée des incidents majeurs liés aux TIC} (Chapitre III, Article 19 de DORA).

DORA impose aux entités financières de reporter à leur autorité compétente (l'ACPR en France) les incidents majeurs selon une taxonomie standardisée (causes, vecteurs, temps de résolution, pertes financières directes et indirectes). À terme, le règlement prévoit la centralisation de ces rapports au sein d'un pôle de l'Union européenne (EU Hub).

La constitution de ce registre européen résoudra structurellement les trois failles actuelles des bases de données :
\begin{itemize}
    \item \textbf{Suppression du biais de déclaration :} Le \textit{reporting} étant réglementairement obligatoire, la base sera exhaustive.
    \item \textbf{Suppression du biais de sélection :} Le registre couvrira toutes les strates du marché financier (grandes banques, mais aussi courtiers, mutuelles et sociétés de gestion), offrant enfin une granularité suffisante pour conditionner la sévérité par taille d'entité.
    \item \textbf{Visibilité sur les tiers :} Le registre DORA cartographiera l'implication des prestataires tiers ICT, validant empiriquement les modèles de contagion externe via la copule de Gumbel présentés en Partie 3.
\end{itemize}

Dès lors, l'architecture du Modèle Interne Partiel développée dans ce mémoire (LDA, théorie des valeurs extrêmes, facteurs de Vasicek, allocation d'Euler) prend toute sa valeur. Le modèle mathématique et les scripts de simulation construits sont intégralement agnostiques à la source de données. Lorsqu'une base d'incidents DORA harmonisée sera rendue disponible à la profession actuarielle, il suffira de substituer ces nouvelles données aux actuelles bases PRC et OpRisk. Le moteur de calcul absorbera ces entrées qualifiées pour produire un SCR cyber d'une précision inédite, répondant simultanément aux exigences qualitatives de DORA et aux contraintes quantitatives de Solvabilité II.

\paragraph{Automatisation du reclassement par NLP.} En attendant ce registre harmonisé, une partie du travail de fiabilisation des données peut d'ores et déjà être outillée. Le reclassement manuel par grille de mots-clés opéré en~2.5 sur la taxonomie Hackmageddon (nécessaire pour isoler le ransomware, absent en tant que catégorie en 2023) illustre une limite récurrente des bases de sinistralité cyber : une taxonomie déclarative instable, dont la correction repose aujourd'hui sur un travail manuel documenté mais non reproductible à grande échelle. Les travaux R\&D menés chez Nexialog Consulting sur l'exploitation du NLP et des modèles de langage pour l'enrichissement de bases de sinistralité cyber (rapports d'incidents, bases de vulnérabilités, publications techniques) ouvrent une voie directe pour systématiser ce type de reclassement --- en le rendant à la fois plus rigoureux (classification supervisée validée statistiquement plutôt que grille de mots-clés ad hoc) et transposable à d'autres bases (OpRisk, futur registre DORA). Cette piste n'a pas été mise en oeuvre dans le présent mémoire, faute de corpus annoté disponible dans le périmètre de ce travail, mais constitue une perspective directe de prolongement.

%--------------------------------------------------------------
\subsection{La PRC comme proxy de fréquence : périmètre, biais, justification}
%--------------------------------------------------------------

La base \emph{Privacy Rights Clearinghouse} (PRC) constitue le socle de
calibration de la fréquence annuelle de référence $\lambda_{ref}$ utilisée
dans le modèle. Son intérêt principal tient à son périmètre défini et à sa
couverture longitudinale, qui permettent d'obtenir un niveau de fréquence
directement exploitable pour une simulation actuarielle.

La PRC joue toutefois un second rôle, plus fragile : celui de source de
sévérité indirecte, via une conversion dite de Jacobs, qui traduit un nombre
d'enregistrements compromis en un montant de perte financière à partir d'un
modèle log-log. Ce mécanisme permet d'élargir le spectre de tailles
d'incidents disponibles, mais introduit un biais de modèle supplémentaire :
la sévérité qui en résulte n'est pas observée directement, elle est
\emph{dérivée}. La calibration en queue de distribution obtenue sur la PRC
--- réalisée en propre sur les 15\,053 incidents 2019--2025 de la base brute,
via la conversion Jacobs $\ln(L) = 7{,}68 + 0{,}76\,\ln(X)$ --- donne un
paramètre de forme $\hat{\xi}_{\text{PRC}} = 1{,}03$ (IC90\%
$[0{,}97~;~1{,}09]$), caractéristique d'une queue extrêmement lourde, à
espérance infinie --- un résultat qui, pris seul, rendrait toute simulation
numériquement instable et qui doit donc être discuté conjointement avec les
corrections méthodologiques présentées en 2.4.

La justification du recours à la PRC malgré ces limites tient à l'absence
d'alternative comparable pour la fréquence : aucune des deux autres sources
mobilisées (Hackmageddon, OpRisk Global) ne permet, seule, d'ancrer un niveau
annuel de sinistralité directement transposable à un portefeuille assurantiel.

%--------------------------------------------------------------
\subsection{Hackmageddon : apports et limites}
%--------------------------------------------------------------

\paragraph{Périmètre consolidé.} La base Hackmageddon a été consolidée sur le
semestre complet et contigu janvier--juin 2026, soit \textbf{1\,041
incidents}. La couverture pluriannuelle initialement envisagée a été
abandonnée : l'année 2024 est absente du jeu transmis, et 2025 n'est
disponible que partiellement. La base est donc exploitée pour la
\emph{structure} du risque, non pour son \emph{niveau} ou sa dynamique
temporelle --- ce qui explique pourquoi elle n'intervient pas dans la
calibration de $\lambda_{ref}$, réservée à la PRC.

\paragraph{Apport 1 : structure du risque et mapping DORA.} Sur les 1\,041
incidents, 840 disposent d'un vecteur d'accès initial identifié (81~\%),
regroupés en familles et reliés directement aux articles DORA :

\begin{center}
\renewcommand{\arraystretch}{1.3}
\begin{tabular}{>{\raggedright}p{5cm} c >{\raggedright\arraybackslash}p{6cm}}
\toprule
\rowcolor{lightblue}
\textbf{Vecteur d'attaque} & \textbf{Part} & \textbf{Exigence DORA} \\
\midrule
Phishing / Ingénierie sociale & 38{,}8~\% & Art.~13 --- Sensibilisation \& formation \\
Exploitation vuln.\ exposée & 33{,}8~\% & Art.~26 --- Tests TLPT ; Art.~8--10 \\
Supply chain / Tiers ICT & 15{,}8~\% & Art.~28--44 --- Gestion risque tiers \\
Identifiants compromis & 6{,}3~\% & Art.~9 --- Contrôle d'accès \& MFA \\
\bottomrule
\end{tabular}
\end{center}

Deux chiffres sont structurants pour la construction des scénarios de la
Partie~4 : \textbf{15{,}8~\%} des attaques empruntent le vecteur tiers/supply
chain, coeur des articles~28--44 ; et \textbf{49{,}6~\%} relèvent de vecteurs
qu'un test TLPT (art.~26) est censé détecter en amont (exploitation et
tiers cumulés).

\paragraph{Limite 1 : biais de visibilité.} Hackmageddon recense des
incidents rendus publics, majoritairement médiatiques ou déclarés, ce qui
sur-représente structurellement les événements les plus visibles au détriment
des incidents ordinaires, moins commentés. La base ne peut donc servir à
estimer un niveau absolu de fréquence, seulement une répartition relative
entre vecteurs d'attaque.

\paragraph{Limite 2 : taxonomie instable.} Une comparaison de structure entre
2023 (3\,019 incidents, Q3 manquant) et 2026 a mis en évidence un artefact de
taxonomie. La catégorie \og Ransomware\fg{} n'existait pas en 2023, les incidents
correspondants étant alors classés sous \og Malware\fg{}. Une lecture naïve
suggérait une hausse du ransomware de +11 points ; après reclassement par
grille de mots-clés commune, le ransomware \emph{recule} en réalité de
23 points.

\begin{center}
\renewcommand{\arraystretch}{1.3}
\begin{tabular}{>{\raggedright}p{4.5cm} c c c}
\toprule
\rowcolor{lightblue}
\textbf{Dimension} & \textbf{2023} & \textbf{2026} & \textbf{$\Delta$} \\
\midrule
Cybercriminalité (motivation) & 81{,}1~\% & 73{,}9~\% & $-7{,}1$ \\
Cyberespionnage (motivation) & 9{,}6~\% & 20{,}6~\% & $+11{,}0$ \\
Ransomware (reclassé) & 35{,}8~\% & 12{,}7~\% & $-23{,}1$ \\
\bottomrule
\end{tabular}
\end{center}

Seules les \textbf{motivations} (taxonomie stable de haut niveau) sont donc
comparables directement dans le temps ; les \textbf{types d'attaque} ne le
sont qu'après un travail de reclassement documenté. Le signal robuste qui
demeure est la bascule du cybercrime de masse vers l'espionnage ciblé, une
information pertinente pour la discussion du risque systémique en Partie~5,
mais qui illustre aussi la fragilité méthodologique de toute lecture brute
d'une base de ce type.

%--------------------------------------------------------------
\subsection{SAS OpRisk Global : proxy de sévérité ICT, lien avec l'approche AMA de Bâle 3}
%--------------------------------------------------------------

\paragraph{Nature et périmètre.} OpRisk Global est la base de référence
historique du risque opérationnel au sens de Bâle~2/3, comptant
\textbf{41\,975 incidents} après nettoyage d'une aberration de saisie. Le
risque cyber n'y constitue pas une catégorie native : il est reconstruit à
partir des catégories \og Systems Security\fg{} et \og Business Disruption and
System Failures\fg{}. Le périmètre cyber $\times$ finance retenu compte
\textbf{582 incidents} sur la période 2000--2026.

\paragraph{Lien avec l'approche AMA.} Cette base est historiquement celle
qui alimentait les modèles internes de type \emph{Advanced Measurement
Approach} évoqués en Partie~1, avant leur retrait au profit de la SMA sous
Bâle~III/IV. Son usage ici renoue donc, à des fins de modélisation interne et
non de calcul réglementaire de Pilier~1, avec la logique originelle de la LDA
bancaire : calibrer une sévérité extrême à partir de pertes réellement
observées plutôt qu'à partir d'un indicateur de volume forfaitaire.

\paragraph{Statistiques de sévérité.}

\begin{center}
\renewcommand{\arraystretch}{1.3}
\begin{tabular}{l r l r}
\toprule
\rowcolor{lightblue}
\textbf{Indicateur} & \textbf{Valeur} & \textbf{Indicateur} & \textbf{Valeur} \\
\midrule
Médiane & 1{,}50 M\$ & P99 & 402 M\$ \\
Moyenne & 26{,}4 M\$ & Maximum & 1\,500 M\$ \\
P90 & 51{,}5 M\$ & Total cumulé & 15{,}4 Md\$ \\
\bottomrule
\end{tabular}
\end{center}

La concentration extrême de la sévérité --- les 10~\% d'incidents les plus
graves représentent 82{,}7~\% de la perte totale, et le seul top~1~\% en
représente 31~\% --- constitue une signature de queue lourde qui justifie
directement le recours à la théorie des valeurs extrêmes développée en
Partie~3.

\paragraph{Choix du seuil POT.} Le choix du seuil $u$ au-dessus duquel la loi
de Pareto généralisée est calibrée résulte d'un compromis classique
biais-variance, arbitré graphiquement. Le \emph{mean excess plot}
(figure~\ref{fig:mep}) présente une pente positive et globalement linéaire
au-delà de $u = 20$~M\euro, signature attendue d'un régime GPD à queue lourde
($\xi > 0$) ; c'est ce seuil, correspondant au percentile~85 de la
distribution des montants, qui est retenu pour la suite.

\begin{figure}[htbp]
\centering
\includegraphics[width=0.85\textwidth]{figures/mean_excess_oprisk.png}
\caption{Mean excess plot sur le périmètre cyber$\times$finance d'OpRisk
Global. La linéarité de $e(u)$ au-delà du seuil retenu $u = 20$~M\euro
(pointillés) confirme le régime de queue lourde justifiant l'ajustement GPD.}
\label{fig:mep}
\end{figure}

\paragraph{Calibration GPD et validation croisée.} Au seuil du percentile~85
($u = 20{,}0$~M\euro, 91 excès), la calibration GPD donne :
\[
\hat{\xi} = 0{,}595 \quad (\text{IC}_{90\%}=[0{,}31\,;\,0{,}83]), \qquad
\hat{\sigma} = 58{,}0~\text{M\euro}, \qquad
\text{VaR}_{99{,}5\%} = 663~\text{M\euro}.
\]
L'intervalle de confiance bootstrap sur cette VaR est $[411\,;\,1037]$~M\euro,
soit un facteur~2{,}5 entre les bornes --- une ampleur qui, à elle seule,
illustre que le SCR ne peut être présenté comme un point.

La comparaison avec la PRC ($\hat{\xi}_{\text{PRC}} = 1{,}03$ contre
$\hat{\xi}_{\text{OpRisk}} = 0{,}595$) révèle un écart de régime de queue,
expliqué par le biais de taille d'OpRisk : les grandes entités financières y
sont sur-représentées et absorbent mieux les chocs, d'où une queue
apparemment plus légère. Les deux sources confirment néanmoins le même régime
qualitatif de queue lourde ($\xi > 0$). Il n'existe donc pas de hiérarchie
\og source primaire\fg{} en sévérité : OpRisk fournit des montants réels, la PRC
un spectre de tailles plus large via la conversion Jacobs --- deux
estimations indépendantes et complémentaires, dont l'écart constitue
lui-même une information sur l'incertitude du modèle.

La décomposition par sous-catégorie révèle en outre deux profils de risque
distincts : les pannes (\og Business Disruption\fg{}) sont moins fréquentes
(19~\% des incidents) mais nettement plus coûteuses, avec une médiane cinq
fois supérieure à celle des intrusions (\og Systems Security\fg{}), suggérant que
les multiplicateurs DORA pourraient à terme être différenciés non seulement
en fréquence mais aussi en sévérité --- sous réserve de la taille réduite de
l'échantillon \og pannes\fg{} (28 excès), qui rend cette calibration indicative.

%--------------------------------------------------------------
\subsection{Cadre épistémique : ce qu'on sait faire et ce qu'on ne sait pas faire}
%--------------------------------------------------------------

La discipline méthodologique retenue dans ce mémoire consiste à ne jamais
prétendre inférer davantage que ce que les données permettent réellement, ce
qui impose d'expliciter aussi bien les capacités que les limites du dispositif.

\paragraph{Ce que l'on sait faire.} On sait calibrer une fréquence de
référence annuelle à partir d'un proxy sectoriel documenté (PRC), et la
moduler selon une structure d'attaque observée empiriquement (Hackmageddon).
On sait calibrer une sévérité de queue lourde par la théorie des valeurs
extrêmes à partir de pertes réellement observées (OpRisk Global), et
confronter cette calibration à une source indépendante pour objectiver
l'incertitude. On sait enfin traiter l'écart entre sources non comme un bruit
à éliminer, mais comme une dimension d'incertitude à part entière, qui
s'ajoute à l'incertitude de paramètre estimée par bootstrap.

\paragraph{Ce que l'on ne sait pas faire.} On ne sait pas mesurer directement
le niveau de conformité DORA d'une entité donnée ; c'est précisément la
raison d'être de la variable latente introduite en 1.4. On ne sait pas
disposer d'une base d'incidents harmonisée à l'échelle européenne, exhaustive
et non biaisée par la visibilité médiatique. On ne sait pas non plus estimer
avec précision la corrélation entre les différentes briques de perte,
faute d'historique suffisant sur les événements joints ; cette incertitude
est traitée en Partie~3 par une structure de dépendance explicite plutôt que
par une hypothèse d'indépendance non fondée.

\paragraph{Deux corrections méthodologiques transparentes.} Ce cadre
épistémique a conduit à identifier et corriger deux défauts du moteur de
simulation en cours de travail. D'une part, avec $\hat{\xi}_{\text{PRC}} =
1{,}03 \geq 1$, l'espérance de la sévérité PRC est infinie et la VaR simulée
diverge numériquement ; un plafond de sévérité de 40~M\euro, ancré sur la
capacité de réassurance cyber observable sur le marché, est appliqué à la
\emph{seule source PRC}. La calibration OpRisk, dont l'indice de queue
$\hat{\xi}_{\text{OpRisk}} = 0{,}595 < 1$ garantit une espérance finie, n'est
en revanche pas plafonnée : ses montants étant observés de bout en bout, y
appliquer un plafond tronquerait des pertes réelles et contredirait sa propre
calibration. D'autre part, le moteur attribuait initialement une sévérité de
queue (GPD) à l'intégralité des événements simulés, au lieu de la seule
fraction $p_u$ dépassant le seuil de queue ; cette correction ramène les SCR
simulés dans des ordres de grandeur cohérents avec les données observées.

La conséquence directe de ce cadre épistémique est méthodologique autant que
présentationnelle : le résultat final ne peut être livré comme un montant
unique, mais doit être présenté comme une distribution large, conditionnelle
à des hypothèses documentées et à des sources dont les écarts sont eux-mêmes
informatifs --- ce principe structure l'ensemble de la Partie~5.

%--------------------------------------------------------------
\subsection{Suffisance du dispositif de données pour l'exercice actuariel}
%--------------------------------------------------------------

Le dispositif retenu --- PRC pour la fréquence, OpRisk Global pour la
sévérité observée, Hackmageddon pour la structure du risque --- appelle une
question préalable à toute discussion de robustesse : ce triptyque est-il
suffisant pour l'exercice mené dans ce mémoire ? La réponse tient à la nature
même de l'exercice. L'objectif n'est pas de produire une formule
réglementaire de Pilier~1, qui exigerait une base exhaustive et représentative
du marché au sens de l'Article~19 --- base qui, comme l'a établi la
section~2.1, n'existe aujourd'hui pour aucun acteur, public ou privé, sur le
risque cyber. L'objectif est de construire, dans l'esprit du Pilier~2 et de
l'ORSA, une mesure prospective dont la valeur ajoutée actuarielle réside dans
la transparence de ses limites autant que dans la précision de son résultat.
Sous cet angle, les trois sources mobilisées couvrent exactement les trois
rôles qu'exige une architecture fréquence-sévérité-structure, et leurs biais
respectifs --- sélection, déclaration, survie, visibilité, instabilité
taxonomique --- sont documentés plutôt que dissimulés (section~2.2).

\paragraph{Pourquoi ne pas simplement en ajouter d'autres.} Des bases
supplémentaires existent (bases commerciales de sinistres cyber assurantiels,
enquêtes sectorielles européennes) et pourraient, en principe, enrichir le
dispositif. Mais chacune porterait ses propres biais de sélection et de
représentativité, sans lever la limite structurelle identifiée dès la
section~2.1 : l'absence, à ce jour, d'un registre harmonisé et exhaustif à
l'échelle du marché. Multiplier les sources déplacerait la difficulté sans la
résoudre, et complexifierait un narratif dont la clarté --- trois sources, trois
rôles distincts, aucune hiérarchie de \og source primaire\fg{} --- constitue
en elle-même un gage de rigueur méthodologique. L'analyse de sensibilité menée
en~5.4 confirme d'ailleurs, \emph{a posteriori}, la pertinence de ce choix :
c'est l'incertitude sur l'indice de queue $\xi$ de la sévérité --- estimé sur
les sources déjà mobilisées --- qui domine très largement toute autre source
de variation du modèle, multipliers de fréquence et paramètres de dépendance
compris. Approfondir l'exploitation des trois sources retenues (bootstrap,
validation croisée, diagnostics de stabilité) réduit donc l'incertitude du
capital de façon bien plus significative que l'ajout d'une quatrième base de
sinistralité.

%--------------------------------------------------------------
\subsection{Validation croisée externe : ENISA Threat Landscape}
%--------------------------------------------------------------

La structure du risque tirée de Hackmageddon (section~2.6), bien que
documentée quant à ses limites de visibilité et de taxonomie, gagne à être
confrontée à une source indépendante, de nature différente et couvrant un
périmètre réglementaire proche de celui du présent mémoire. Le rapport
\emph{ENISA Threat Landscape 2025}, déjà mobilisé en~4.3 pour ancrer les
bornes du multiplicateur associé à l'exploitation de vulnérabilités, remplit
ce rôle : constitué par l'agence européenne de cybersécurité elle-même à
partir de 4\,875 incidents de périmètre européen, il repose sur une
méthodologie de collecte distincte de celle, essentiellement médiatique, de
Hackmageddon.

La confrontation porte sur un point précis, déjà exploité en~4.3 : ENISA
mesure un taux de conversion tentative-vers-intrusion de 70~\% pour
l'exploitation de vulnérabilités, contre seulement 27~\% pour le phishing.
Cette asymétrie, mesurée par un organisme réglementaire sur un périmètre
indépendant, confirme un résultat qu'Hackmageddon, par construction, ne peut
révéler à partir du seul volume d'incidents observés : le phishing domine la
structure en volume (38{,}8~\% contre 33{,}8~\% pour l'exploitation de
vulnérabilités, section~2.6), mais convertit nettement moins efficacement en
intrusion réussie. Les deux sources s'accordent ainsi, chacune sur son propre
plan --- structure des incidents rapportés pour Hackmageddon, taux de
conversion pour ENISA --- sur une même hiérarchie qualitative des vecteurs
d'attaque, ce qui renforce la crédibilité du mapping vecteurs~$\rightarrow$~exigences
DORA retenu en~4.2 sans nécessiter de recalibration.

\paragraph{Piste non mobilisée : Verizon DBIR.} Le \emph{Data Breach
Investigations Report} de Verizon, fondé sur des méthodologies de collecte
distinctes à la fois de Hackmageddon et d'ENISA, constituerait une troisième
confrontation naturelle. Il n'est pas mobilisé dans le présent mémoire, faute
de statistiques suffisamment vérifiées dans le périmètre de ce travail pour
en garantir l'exactitude des chiffres cités --- une prudence identique à
celle déjà retenue pour les références bibliographiques marquées comme à
vérifier (cf.~\texttt{references.bib}). Cette confrontation constitue une
perspective directe de prolongement, dans le même esprit que l'automatisation
du reclassement par NLP déjà évoquée en~2.4.

%==============================================================
\section{Modélisation}
%==============================================================

Les deux parties précédentes ont posé le cadre réglementaire et les données
disponibles. Cette troisième partie formalise le moteur actuariel qui
transforme ces éléments en une distribution de pertes, puis en une mesure de
capital. L'architecture retenue reste fondée sur une approche par
distribution des pertes (LDA), enrichie pour tenir compte de la
multiplicité des sources de risque et de leur dépendance mutuelle.

%--------------------------------------------------------------
\subsection{Architecture LDA : les quatre briques de perte}
%--------------------------------------------------------------

L'approche LDA classique calibre une fréquence de sinistres stationnaire,
indépendante du niveau de maîtrise opérationnelle de l'assureur. Cette
hypothèse est précisément ce que le règlement DORA invite à remettre en
cause : un défaut de conformité aux exigences de résilience numérique
n'affecte pas la sévérité des sinistres de façon uniforme, mais leur
\emph{fréquence}, par l'intermédiaire des vecteurs d'attaque que chaque
article de DORA est censé atténuer.

Le modèle retenu décompose la perte physique simulée $L_{DORA}$ en trois briques, chacune dotée de sa propre logique :
\begin{equation}
L_{DORA} = L_{\text{remédiation}} + L_{\text{sanction}} + L_{\text{prestataire}}
\end{equation}

La brique d'\textbf{aggravation} ne constitue pas une perte simulée supplémentaire s'ajoutant à cette somme, mais une mesure dérivée. Elle quantifie le surcroît de capital entre un état conforme et un état non conforme, au travers du différentiel $\Delta_{DORA} = VaR_\alpha(L_{DORA}^{\text{non conforme}}) - VaR_\alpha(L_{DORA}^{\text{conforme}})$, développé en Partie 4.

La brique de \textbf{remédiation} constitue le coeur fréquence-sévérité du
modèle : elle représente le coût direct de gestion des incidents cyber,
simulé par composition d'une fréquence (section~3.2) et d'une sévérité
extrême (section~3.3). La brique de \textbf{sanction} traduit le risque
réglementaire explicite prévu par les articles 50-52 de DORA, pilotée par la
probabilité de défaut de conformité $PCD_i(\theta)$ introduite en~1.4. La
brique de \textbf{prestataire} modélise un événement indépendant impliquant
un tiers ICT critique, dont la probabilité d'occurrence diminue lorsque
l'entité maîtrise mieux la gestion de ses tiers (chapitre~V de DORA). Enfin,
la brique d'\textbf{aggravation} ne constitue pas une perte simulée
supplémentaire au même titre que les trois précédentes : elle mesure le
surcroît de capital entre un état conforme et un état non conforme, au
travers de la différence $\Delta_{DORA} = VaR_\alpha(L_{DORA}^{\text{non
conforme}}) - VaR_\alpha(L_{DORA}^{\text{conforme}})$, développée en
Partie~4.

Cette architecture rompt avec une logique de chargements multiplicatifs
appliqués uniformément sur une sévérité globale : chaque brique possède sa
propre distribution, sa propre logique métier, et son propre lien --- ou
absence de lien --- à la variable latente de conformité. Elle impose en
contrepartie de traiter explicitement la question de la dépendance entre
briques (section~3.4), qu'une architecture additive naïve négligerait.

%--------------------------------------------------------------
\subsection{Fréquence : NegBin, calibration PRC, sur-dispersion}
%--------------------------------------------------------------

\paragraph{Décomposition par vecteur d'attaque.} La fréquence de référence
$\lambda_{ref}$, calibrée sur la base PRC ($\lambda_{ref} = 341$
incidents/an), est décomposée en sous-fréquences par vecteur d'attaque,
pondérées par la structure observée dans la base Hackmageddon :
\begin{equation}
\lambda_{ref} = \sum_{v} \lambda_v,
\qquad
\lambda_v = \lambda_{ref} \times \pi_v,
\end{equation}
où $\pi_v$ est la part du vecteur $v$ dans la structure Hackmageddon
(38{,}8~\% phishing, 33{,}8~\% exploitation de vulnérabilités, 15{,}8~\%
supply chain, 6{,}3~\% identifiants compromis, 5{,}3~\% autres). Cette
décomposition transforme un constat structurel en levier de modélisation :
un défaut de conformité sur une exigence donnée n'agit que sur la
sous-fréquence du vecteur qu'elle est censée mitiger, pas sur l'ensemble de
$\lambda_{ref}$.

\paragraph{Choix de la loi binomiale négative.} La fréquence agrégée
$\lambda_{S_x} = \sum_v \lambda_{ref}\,\pi_v\, m_{v,S_x}$ obtenue sous un
scénario $S_x$ n'est pas simulée par une loi de Poisson, mais par une loi
binomiale négative, pour tenir compte de la sur-dispersion observée dans les
données ($\sigma^2/\mu > 1$). Le facteur de sur-dispersion
$\phi = \sigma^2/\mu = 9{,}2$ est traité comme une propriété structurelle de
la sinistralité, supposée indépendante du niveau de fréquence : c'est donc
$\phi$ qui est maintenu constant d'un scénario à l'autre, et non le
paramètre $r$ de la binomiale négative, qui est au contraire recalculé sous
chaque scénario :
\begin{equation}
r_{S_x} = \frac{\lambda_{S_x}}{\phi - 1},
\qquad
p_{S_x} = \frac{r_{S_x}}{r_{S_x} + \lambda_{S_x}}.
\end{equation}
Augmenter $\lambda$ via les multiplicateurs DORA déplace ainsi la fréquence
moyenne tout en préservant la \emph{forme} de sur-dispersion de la
sinistralité --- une hypothèse de modélisation explicite, dont la limite est
que le comportement de la sur-dispersion en régime de forte dégradation
n'est pas lui-même observé, faute d'historique de crise DORA.

\paragraph{De la fréquence par scénario à la fréquence par entité.} Le choix
d'un scénario $S_x \in \{S_0, S_1, S_2\}$ reste, par construction, discret et
appliqué uniformément à l'ensemble du marché. La variable latente de
conformité $PCD_i(\theta)$ introduite en~1.4 permet de dépasser cette
limite, en substituant à ce choix discret une interpolation continue :
\begin{equation}
m_v\bigl(PCD_i(\theta)\bigr) = 1 + PCD_i(\theta) \times \left(\bar{m}_v^{S_2} - 1\right),
\end{equation}
où $\bar{m}_v^{S_2}$ est le centre de l'intervalle de multiplicateur sourcé
en scénario $S_2$. Une entité au profil de conformité mature conserve ainsi
une fréquence proche de $S_0$ même en cas de choc systémique $\theta$
défavorable, tandis qu'une entité en retard de conformité peut, sous un
$\rho_i$ élevé, dépasser le multiplicateur $S_2$ dès l'environnement normal.

%--------------------------------------------------------------
\subsection{Sévérité : EVT/GPD, calibration OpRisk, validation croisée PRC}
%--------------------------------------------------------------

\paragraph{Choix de la loi de Pareto généralisée.} La sévérité des sinistres
cyber présente une signature de queue lourde documentée en Partie~2 : sur
OpRisk Global, les 10~\% d'incidents les plus graves concentrent 82{,}7~\%
de la perte totale. Cette concentration extrême justifie le recours à la
théorie des valeurs extrêmes, et plus précisément à une modélisation de la
queue de distribution par une loi de Pareto généralisée (GPD) au-delà d'un
seuil $u$, selon l'approche \emph{peaks-over-threshold}.

\paragraph{Calibration OpRisk.} Au seuil du percentile~85
($u = 20{,}0$~M\euro, 91 excès), la calibration donne :
\begin{equation}
\hat{\xi}_{\text{OpRisk}} = 0{,}595 \quad (\text{IC}_{90\%}=[0{,}31\,;\,0{,}83]),
\qquad
\hat{\sigma}_{\text{OpRisk}} = 58{,}0~\text{M\euro},
\qquad
\text{VaR}_{99{,}5\%} = 663~\text{M\euro}.
\end{equation}
L'intervalle de confiance bootstrap sur cette VaR, $[411\,;\,1037]$~M\euro,
matérialise à lui seul un facteur~2{,}5 entre les bornes, uniquement
imputable à l'incertitude de paramètre.

\paragraph{Stabilité du paramètre de forme.} Le choix du seuil est corroboré
par deux diagnostics complémentaires. Le \emph{diagnostic GPD}
(figure~\ref{fig:diagnostic-gpd}) montre que $\hat{\xi}$ se stabilise autour
de $0{,}6$ dès que le seuil atteint $u = 20$~M\euro, tout en conservant un
nombre d'excès suffisant ($n = 91 > 30$) pour une estimation fiable : le
seuil retenu se situe précisément à la frontière où le paramètre de forme
quitte la zone instable ($\xi \geq 1$) pour se fixer dans un régime de
moments finis. Le \emph{Hill plot} (figure~\ref{fig:hill-ic}) confirme
qualitativement ce niveau, l'estimateur de Hill restant compatible avec la
référence MLE $\hat{\xi} \approx 0{,}60$ dans la plage des extrêmes retenus.

\begin{figure}[htbp]
\centering
\includegraphics[width=\textwidth]{figures/diagnostic_gpd_oprisk.png}
\caption{Diagnostic GPD sur le périmètre cyber$\times$finance d'OpRisk.
\textit{Gauche} : stabilité de $\hat{\xi}$ en fonction du seuil $u$ ; le
paramètre passe sous $\xi = 1$ (espérance finie) et se stabilise autour de
$0{,}6$ au seuil retenu $u = 20$~M\euro. \textit{Droite} : le nombre d'excès
correspondant ($n = 91$) demeure au-dessus du minimum de $30$ requis pour une
estimation robuste.}
\label{fig:diagnostic-gpd}
\end{figure}

\paragraph{L'estimateur de Hill : définition et rôle diagnostique.}
Introduit par \citet{Hill1975}, l'estimateur de Hill fournit une estimation
de l'indice de queue $\xi$ à partir des $k$ statistiques d'ordre les plus
élevées d'un échantillon $X_{(1)} \leq \dots \leq X_{(n)}$, sans passer par
l'ajustement paramétrique d'une loi de Pareto généralisée :
\begin{equation}
\hat{\xi}_{\text{Hill}}(k) = \frac{1}{k}\sum_{i=1}^{k}
\ln X_{(n-i+1)} - \ln X_{(n-k)}.
\label{eq:hill}
\end{equation}
Contrairement au maximum de vraisemblance sur la GPD, qui exploite le cadre
théorique de Pickands--Balkema--de~Haan et estime conjointement $\xi$ et
$\sigma$, l'estimateur de Hill repose sur une hypothèse de Pareto \emph{pure}
dans la queue et ne fournit qu'un indice de forme, sans paramètre d'échelle
explicite \citep{Beirlant2004}. Sa propriété la plus documentée dans la
littérature EVT est un compromis biais--variance piloté par le choix de
$k$~: pour $k$ trop petit, l'estimateur est instable (variance élevée,
intervalle de confiance large) ; pour $k$ trop grand, il intègre des
observations situées hors du régime purement asymptotique de la queue, ce
qui induit un biais positif croissant avec $k$ \citep{Embrechts1997,
Beirlant2004}.

\paragraph{Application et confrontation au MLE-GPD.} Le Hill plot obtenu sur
le périmètre cyber$\times$finance d'OpRisk (figure~\ref{fig:hill-ic})
illustre ce compromis de façon nette~: l'estimateur croît de façon quasi
monotone avec $k$, passant de $\hat{\xi} \approx 0{,}5$ pour $k$ faible à
$\hat{\xi} \approx 1{,}8$ pour $k = 200$, tandis que la bande de confiance
asymptotique à 90~\% s'élargit continûment. La référence par maximum de
vraisemblance ($\hat{\xi}_{\text{MLE}} = 0{,}595$, pointillés rouges) n'est
statistiquement compatible avec l'estimateur de Hill que sur une plage
restreinte, $k \in [3, 25]$ environ, où la bande de confiance recouvre
encore la valeur MLE. Au-delà de $k \approx 30$, l'estimateur de Hill
diverge continûment de cette référence, ce qui traduit le biais positif
attendu lorsque des observations hors du régime purement asymptotique de la
queue sont intégrées au calcul. Sur le seuil effectivement retenu pour la
calibration ($k = 86$), cet écart atteint 138{,}7~\%
($\hat{\xi}_{\text{Hill}} = 1{,}42$ contre $\hat{\xi}_{\text{MLE}} = 0{,}595$),
ce qui confirme que l'estimateur de Hill n'est pas adapté à ce niveau de
$k$ et que le MLE-GPD, qui exploite un cadre théorique plus général
(loi GPD complète plutôt que Pareto pure), demeure l'estimateur de
référence pour la calibration finale.

\begin{figure}[htbp]
\centering
\includegraphics[width=0.85\textwidth]{figures/hill_plot_ic90_oprisk.png}
\caption{Hill plot avec intervalle de confiance asymptotique à 90~\% sur
OpRisk Global. La référence MLE ($\hat{\xi} \approx 0{,}595$, pointillés
rouges) n'est statistiquement compatible avec l'estimateur de Hill que pour
$k \lesssim 25$~; au-delà, la dérive de $\hat{\xi}_{\text{Hill}}$ traduit le
biais croissant lié à l'intégration d'observations hors du régime
asymptotique pur de la queue.}
\label{fig:hill-ic}
\end{figure}



\paragraph{Validation croisée PRC.} La calibration sur la PRC, via la
conversion Jacobs, donne $\hat{\xi}_{\text{PRC}} = 1{,}03$ (seuil
$u = 4{,}18$~M\euro, 2\,258 excès), supérieur à
$\hat{\xi}_{\text{OpRisk}} = 0{,}595$, les IC90\% des deux estimateurs ne se
chevauchant pas. Ce paramètre correspond formellement à une espérance de
sévérité infinie ($\xi \geq 1$), ce qui impose un plafond de sévérité
opérationnel --- fixé à 40~M\euro, ancré sur la capacité de réassurance cyber
observable sur le marché --- pour que la simulation reste numériquement
stable ; ce plafond ne s'applique qu'à la source PRC, la calibration OpRisk
($\xi < 1$, espérance finie) demeurant non plafonnée. L'écart entre
les deux calibrations s'explique par le biais de taille d'OpRisk : les
grandes entités financières y sont sur-représentées et absorbent mieux les
chocs, d'où une queue apparemment plus légère. Les deux sources confirment
néanmoins le même régime qualitatif de queue lourde ($\xi > 0$), et aucune
hiérarchie de type \og source primaire\fg{} n'est retenue entre elles --- leur
écart est traité comme une composante de l'incertitude totale du modèle,
non comme un désaccord à trancher.

%--------------------------------------------------------------
\subsection{Dépendance : copule de Gumbel, contagion, modèle à facteur commun}
%--------------------------------------------------------------

\paragraph{Pourquoi l'indépendance est intenable.} Une architecture purement
additive des quatre briques, simulées indépendamment, ignorerait la
discussion de contagion menée en~1.5 : une même dégradation de
l'environnement cyber systémique $\theta$ tend à accroître simultanément la
fréquence de remédiation, la probabilité de sanction (via le $PCD_i(\theta)$)
et l'occurrence d'un événement prestataire. Traiter ces briques comme
indépendantes sous-estimerait mécaniquement les quantiles extrêmes de la
perte agrégée, précisément là où le capital réglementaire est évalué.

\paragraph{Modèle à facteur commun.} La contagion \emph{interne}, décrite
en~1.5, est capturée par la structure factorielle de Vasicek déjà introduite
pour le $PCD_i(\theta)$ : le facteur $\theta$ agit comme un choc commun
affectant simultanément la brique de remédiation (via la fréquence) et la
brique de sanction (via le $PCD_i$). Ce choix évite d'introduire une seconde
structure de dépendance ad hoc : le même $\theta$ qui module la conformité
individuelle module également l'intensité de la sinistralité observée.

\paragraph{Copule de Gumbel pour la dépendance de queue.} La contagion
\emph{externe}, portée par les tiers ICT, est en revanche mal représentée
par une simple corrélation gaussienne : le risque prestataire se matérialise
typiquement de façon conjointe avec les événements les plus extrêmes de
remédiation, lors de chocs systémiques affectant un fournisseur mutualisé.
Cette dépendance de queue positive est représentée par une copule de
Gumbel, dont le paramètre $\eta \geq 1$ contrôle l'intensité de la
dépendance dans les scénarios extrêmes :
\begin{equation}
C_\eta(u,v) = \exp\!\left[-\left((-\ln u)^\eta + (-\ln v)^\eta\right)^{1/\eta}\right],
\end{equation}
appliquée entre les rangs simulés de $L_{\text{remédiation}}$ et
$L_{\text{prestataire}}$. Le cas $\eta = 1$ correspond à l'indépendance ; des
valeurs de $\eta$ croissantes renforcent la probabilité que les deux briques
réalisent simultanément une valeur extrême, conformément à l'intuition
selon laquelle une crise cyber sectorielle majeure implique souvent à la
fois une remédiation interne lourde et une défaillance de tiers.

Cette double structure --- facteur commun de Vasicek pour la contagion
interne, copule de Gumbel pour la contagion externe --- permet de conserver
une lecture économique distincte des deux canaux identifiés en~1.5, plutôt
que de les fondre artificiellement dans une unique matrice de corrélation.

%--------------------------------------------------------------
\subsection{Agrégation Monte Carlo et allocation d'Euler}
%--------------------------------------------------------------

\paragraph{Simulation Monte Carlo.} L'agrégation finale des quatre briques
est réalisée par simulation Monte Carlo. À chaque itération : (i)~un facteur
systémique $\theta$ est tiré, déterminant le $PCD_i(\theta)$ de l'entité et
la fréquence effective $\lambda_{S_x}$ ou $\lambda(PCD_i)$ ; (ii)~la
fréquence de chaque vecteur est simulée par la binomiale négative calibrée
en~3.2 ; (iii)~la sévérité de chaque événement est tirée dans la GPD
calibrée en~3.3, sous la source retenue (OpRisk ou PRC) ; (iv)~la dépendance
entre briques est imposée via la structure factorielle et la copule de
Gumbel décrites en~3.4. Le recours à Monte Carlo n'est pas ici un simple
choix technique : la grandeur recherchée n'est pas une perte moyenne, mais
un quantile extrême de la distribution agrégée, dont la forme analytique
n'existe pas sous une architecture à dépendance non gaussienne.

\paragraph{Allocation d'Euler.} Une fois la perte agrégée simulée, il est
nécessaire de répartir le capital total entre les quatre briques, pour
identifier lesquelles contribuent le plus au risque extrême. L'allocation
d'Euler, cohérente avec la propriété d'homogénéité de la VaR, répartit le
capital total proportionnellement à la contribution marginale de chaque
brique conditionnellement au scénario où la perte agrégée dépasse son
quantile $\alpha$ :
\begin{equation}
AC_k = \mathbb{E}\!\left[L_k \,\middle|\, L_{DORA} = VaR_\alpha(L_{DORA})\right],
\qquad
\sum_{k} AC_k = VaR_\alpha(L_{DORA}),
\end{equation}
pour $k \in \{\text{remédiation}, \text{sanction}, \text{prestataire},
\text{aggravation}\}$. En pratique, cette espérance conditionnelle est
estimée numériquement par une fenêtre de tirages Monte Carlo au voisinage du
quantile $\alpha$ recherché, faute de forme fermée pour l'architecture de
dépendance retenue.

Cette allocation permet de répondre à une question distincte de celle du
niveau global de capital : non pas \og combien faut-il provisionner ?\fg{}, mais
\og quelle brique --- remédiation, sanction, prestataire ou aggravation ---
consomme le plus de capital dans les scénarios extrêmes ?\fg{}. Cette
décomposition est présentée et discutée en Partie~5, où elle est mise en
regard des travaux antérieurs sur la structure du risque cyber.

\vspace{1em}


%==============================================================
\section{Scénarios de conformité DORA}
%==============================================================

La Partie~3 a posé l'architecture de modélisation ; cette quatrième partie
formalise la manière dont le niveau de conformité DORA se traduit
concrètement dans les paramètres du modèle, puis se propage jusqu'à une
mesure de surcoût de capital, le $\Delta_{DORA}$.

%--------------------------------------------------------------
\subsection{Les trois niveaux de non-conformité}
%--------------------------------------------------------------

Le modèle retient trois scénarios-types de conformité, $S_x \in \{S_0, S_1,
S_2\}$, qui bornent l'espace des situations observables sur le marché.
$S_0$ correspond à une entité pleinement conforme aux cinq piliers DORA
détaillés en~1.2 ; $S_2$ correspond à une situation de non-conformité
totale ; $S_1$ représente une conformité partielle, intermédiaire entre les
deux régimes précédents.

Ces trois scénarios ne sont pas de simples repères qualitatifs : ils
ancrent les bornes numériques des multiplicateurs de fréquence calibrés en
4.3, et servent également de cas limites à l'interpolation continue portée
par la variable latente $PCD_i(\theta)$ introduite en~1.4. Le scénario $S_0$
correspond ainsi à la limite $PCD_i(\theta) \to 0$, et $S_2$ à la limite
$PCD_i(\theta) \to 1$, ce qui permet de faire cohabiter une lecture par
scénarios de marché, utile pour la communication des résultats, et une
lecture par profil d'entité, plus fine, mobilisée dans les applications
individuelles du modèle.

%--------------------------------------------------------------
\subsection{Mapping vecteurs d'attaque $\rightarrow$ exigences DORA}
%--------------------------------------------------------------

La fréquence de référence $\lambda_{ref} = 341$ incidents/an, calibrée sur la
PRC, est décomposée en sous-fréquences par vecteur d'attaque, pondérées par
la structure observée dans Hackmageddon (1\,041 incidents, S1~2026) et
reliées explicitement à l'article DORA censé les atténuer :

\begin{center}
\renewcommand{\arraystretch}{1.25}
\begin{tabular}{>{\raggedright}p{4.5cm} c >{\raggedright\arraybackslash}p{6cm}}
\toprule
\rowcolor{lightblue}
\textbf{Vecteur} & \textbf{$\pi_v$} & \textbf{Exigence DORA atténuante} \\
\midrule
Phishing / ingénierie sociale & 38{,}8\,\% & Art.~13 --- Sensibilisation \& formation \\
Exploitation de vulnérabilités & 33{,}8\,\% & Art.~26 --- Tests TLPT \\
Supply chain / tiers ICT & 15{,}8\,\% & Art.~28-44 --- Gestion du risque tiers \\
Identifiants compromis & 6{,}3\,\% & Art.~9 --- Contrôle d'accès et MFA \\
Autres / transverse & 5{,}3\,\% & --- \\
\bottomrule
\end{tabular}
\end{center}

Cette décomposition transforme un constat structurel --- la part de chaque
vecteur dans la sinistralité observée --- en un levier de modélisation
directement actionnable : un défaut de conformité sur une exigence donnée
n'agit que sur la sous-fréquence du vecteur qu'elle est censée mitiger, et
non sur l'ensemble de $\lambda_{ref}$. Cette granularité est ce qui distingue
l'approche retenue d'un simple chargement multiplicatif global appliqué à la
sinistralité totale.

%--------------------------------------------------------------
\subsection{Multiplicateurs de fréquence différenciés par exigence}
%--------------------------------------------------------------

Sous un scénario $S_x$, chaque vecteur $v$ se voit appliquer un
multiplicateur $m_{v,S_x}$, non pas une valeur ponctuelle mais un
\textbf{intervalle} $[m_{min}, m_{max}]$, dont les bornes sont ancrées sur la
littérature empirique relative à l'efficacité des contrôles de sécurité
correspondants.

\paragraph{Phishing (Art.~13).} Ponemon/IBM \emph{Cost of a Data Breach
2025} établit que la formation des employés peut réduire le risque jusqu'à
86~\%, avec un retour sur investissement médian de 5 fois. Cette mesure
étant agrégée sur l'ensemble des incidents et non spécifique au phishing,
nous n'en retenons pas l'inverse littéral, mais une fraction prudente de cet
effet : $[1{,}8\,;\,3{,}0]$ en scénario $S_2$.

\paragraph{Exploitation de vulnérabilités (Art.~26).} ENISA \emph{Threat
Landscape 2025} (4\,875 incidents, périmètre européen) mesure un taux de
conversion tentative-vers-intrusion de 70~\% pour l'exploitation de
vulnérabilités contre 27~\% pour le phishing, soit un ratio directement
observé de $\times 2{,}6$ --- la seule borne de la grille fondée sur une
probabilité de succès conditionnelle mesurée par un organisme réglementaire.
Elle ancre l'intervalle $[2{,}0\,;\,2{,}6]$ en $S_2$.

\paragraph{Supply chain et tiers ICT (Art.~28-44).} IBM/Ponemon mesurent un
surcoût de 11{,}8~\% pour les brèches impliquant un partenaire commercial et
de 8{,}3~\% pour celles impliquant une compromission logicielle ;
SecurityScorecard documente une part des brèches d'origine tierce passée de
15~\% à 36~\% en une année. La combinaison de ce surcoût mesuré et de cette
tendance de marché donne l'intervalle $[1{,}9\,;\,2{,}6]$ en $S_2$.

\paragraph{Identifiants compromis (Art.~9).} Microsoft Research
(arXiv:2305.00945, étude \emph{peer-reviewed} sur comptes Azure Active
Directory réels) établit que l'authentification multifacteur réduit le
risque de compromission de 99{,}22~\% en population générale, et de
98{,}56~\% même en cas de fuite préalable des identifiants. Il s'agit de la
mesure la plus rigoureuse de la grille, mais aussi la plus susceptible
d'être mal transposée à l'échelle d'une entité financière hétérogène : nous
en retenons des bornes tempérées, $[3{,}0\,;\,8{,}0]$ en $S_2$ --- les plus
élevées de la grille, sans reproduire l'inversion littérale que le
99{,}22~\% suggérerait naïvement.

\begin{center}
\renewcommand{\arraystretch}{1.3}
\begin{tabular}{>{\raggedright}p{3.8cm} c >{\centering}p{2.6cm} >{\centering\arraybackslash}p{2.6cm}}
\toprule
\rowcolor{lightblue}
\textbf{Vecteur} & \textbf{$\pi_v$} & \textbf{$S_1$ partiel} & \textbf{$S_2$ non conforme} \\
\midrule
Phishing & 38{,}8\,\% & [1{,}2 ; 1{,}6] & [1{,}8 ; 3{,}0] \\
Exploit. vuln. & 33{,}8\,\% & [1{,}3 ; 1{,}8] & [2{,}0 ; 2{,}6] \\
Supply chain & 15{,}8\,\% & [1{,}3 ; 1{,}9] & [1{,}9 ; 2{,}6] \\
Identifiants & 6{,}3\,\% & [1{,}5 ; 3{,}0] & [3{,}0 ; 8{,}0] \\
\bottomrule
\end{tabular}
\end{center}

\paragraph{Le principe central : l'intervalle, pas le point.} Pour chaque
itération de simulation, un multiplicateur $m_{v,S_x}$ est \textbf{tiré
uniformément} dans son intervalle --- jamais fixé à une valeur unique. Cette
discipline, appliquée dès la calibration des hypothèses et non seulement en
sortie de modèle, permet de propager une incertitude de \emph{calibration},
et non la seule incertitude d'échantillonnage, jusqu'au résultat final.

Sous ce cadre, la fréquence effective sous un scénario $S_x$ s'écrit :
\begin{equation}
\lambda_{S_x} = \sum_{v} \lambda_{ref} \times \pi_v \times m_{v,S_x},
\end{equation}
puis se propage dans la binomiale négative selon le mécanisme de
sur-dispersion constante ($\phi = 9{,}2$) détaillé en~3.2.

%--------------------------------------------------------------
\subsection{Propagation et $\Delta_{DORA}$ : distribution, pas un point}
%--------------------------------------------------------------

Les résultats obtenus, pour $\alpha = 99{,}5~\%$, sous les deux sources de
sévérité, sont les suivants :

\begin{center}
\renewcommand{\arraystretch}{1.3}
\begin{tabular}{>{\raggedright}p{2.5cm} >{\raggedright}p{3.2cm} >{\centering}p{2.8cm} >{\centering\arraybackslash}p{3.8cm}}
\toprule
\rowcolor{lightblue}
\textbf{Source} & \textbf{Scénario} & \textbf{$\Delta_{DORA}$ médiane} & \textbf{IC90\,\%} \\
\midrule
PRC & $S_1$ partiel & 720{,}9 M\euro & [555{,}7~;~895{,}5] M\euro \\
PRC & $S_2$ non conforme & 2\,014{,}7 M\euro & [1\,606{,}9~;~2\,366{,}1] M\euro \\
\addlinespace
OpRisk & $S_1$ partiel & 1\,514{,}4 M\euro & [522{,}5~;~7\,863{,}2] M\euro \\
OpRisk & $S_2$ non conforme & 3\,879{,}3 M\euro & [1\,496{,}9~;~22\,249{,}3] M\euro \\
\bottomrule
\end{tabular}
\end{center}

% NOTE(relecture) : grille régénérée via notebooks/07_bootstrap_delta_dora.py
% (200 tirages x 20 000 sim) APRÈS (a) recalibration de la brique prestataire
% en surcoût relatif et (b) recalibration PRC EN PROPRE sur les données brutes
% (conversion Jacobs ln(L)=7.68+0.76 ln(X), 2258 excès, cf. notebooks/13).
% Les DEUX sources sont désormais bootstrappées sur données réelles (PRC 2258
% excès, OpRisk 91 excès). L'IC PRC reste plus resserré car la conversion
% Jacobs lisse la variance de sévérité (transformation déterministe X->L).
Sous PRC, un défaut de conformité DORA se traduit par une charge de capital
quantifiée et non simplement qualitative : le passage de $S_1$ à $S_2$
accroît le $\Delta_{DORA}$ médian d'un facteur~2{,}8 (720{,}9 à
2\,014{,}7~M\euro), cohérent avec le rapport des multiplicateurs globaux de
fréquence entre les deux scénarios. Ce $\Delta_{DORA}$ mesure le surcoût
strict de capital imputable au défaut de conformité, à profil de risque
constant (contrefactuel à graine figée).

Sous OpRisk, le même défaut de conformité produit un $\Delta_{DORA}$ du même
ordre de grandeur : 1\,514{,}4~M\euro en $S_1$ et 3\,879{,}3~M\euro
en $S_2$, soit un facteur proche de~1{,}9 par rapport à la PRC sur $S_2$. Ce
résultat est notable : une fois la PRC calibrée de bout en bout sur les
données brutes (et non plus sur des paramètres externes non reproduits), les
deux sources indépendantes \emph{convergent} bien davantage qu'anticipé ---
l'écart de niveau entre sources (facteur~1{,}9) devient comparable à l'écart
entre scénarios de conformité sous une même source (facteur~2{,}6 de $S_1$ à
$S_2$ sous OpRisk). Ce qui distingue encore nettement les deux sources n'est
plus le \emph{niveau} du capital mais son \emph{incertitude} : l'IC90\%
OpRisk s'étend sur un facteur~15 entre bornes (bootstrap de seulement 91
excès), là où l'IC PRC, appuyé sur 2\,258 excès et une sévérité lissée par la
conversion Jacobs, reste beaucoup plus resserré (cf.~section suivante).

%--------------------------------------------------------------
\subsection{Bootstrap et intervalles de confiance}
%--------------------------------------------------------------

La distribution du $\Delta_{DORA}$ est obtenue par un bootstrap à deux
niveaux, combiné à une architecture Monte Carlo.

\begin{enumerate}[topsep=4pt]
\item \textbf{Niveau 1 --- incertitude des multiplicateurs.} À chaque
itération, un jeu de multiplicateurs $\{m_{v,S_x}\}$ est tiré dans les
intervalles sourcés en~4.3, produisant un $\lambda_{S_x}$ lui-même
aléatoire plutôt que fixé.
\item \textbf{Niveau 2 --- incertitude de sévérité.} Un couple
$(\hat\xi, \hat\sigma)$ est obtenu par rééchantillonnage bootstrap des excès
réels, \emph{pour les deux sources} : pour OpRisk, sur les 91 excès observés
au-dessus du seuil $u = 20$~M\euro (périmètre cyber$\times$finance de 582
incidents) ; pour la PRC, sur les 2\,258 excès de la sévérité dérivée de la
conversion Jacobs au-dessus du seuil $u = 4{,}18$~M\euro (15\,053 incidents
2019--2025). L'incertitude de sévérité PRC est ainsi elle aussi
\emph{data-driven}, quoique plus resserrée : la conversion Jacobs, en
traduisant de façon déterministe le nombre d'enregistrements en montant,
lisse mécaniquement la variance de sévérité.
\item \textbf{Simulation Monte Carlo.} Pour chaque tirage des niveaux 1 et
2, le modèle LDA est simulé pour obtenir $SCR_{S_0}$ et $SCR_{S_x}$, avec la
\textbf{même graine} de simulation pour les deux régimes, condition
méthodologique centrale garantissant que l'écart observé provient
uniquement du changement de conformité et non de la variance
d'échantillonnage Monte Carlo (principe des \emph{common random numbers}).
\item \textbf{Résultat.} $\Delta_{DORA} = SCR_{S_x} - SCR_{S_0}$ est
accumulé sur l'ensemble des tirages, produisant un intervalle de confiance
plutôt qu'un point.
\end{enumerate}

\paragraph{Asymétrie des intervalles.} L'intervalle de confiance associé à la
source OpRisk est nettement plus large que celui de la source PRC (facteur
proche de~15 contre facteur proche de~1{,}5 sur $S_2$). Cette asymétrie
n'est pas un artefact et --- depuis la recalibration de la PRC sur ses
données brutes --- ne provient plus d'un bootstrap PRC \emph{incomplet} : les
deux sources intègrent désormais l'incertitude de sévérité mesurée par
rééchantillonnage réel (2\,258 excès pour la PRC, 91 pour OpRisk). L'écart
d'amplitude tient à deux causes structurelles. D'une part, le nombre d'excès :
OpRisk ne dispose que de 91 observations de queue, contre 2\,258 pour la PRC,
si bien que le rééchantillonnage OpRisk fait varier $(\hat\xi, \hat\sigma)$
bien davantage. D'autre part, la nature de la sévérité PRC : la conversion
Jacobs, transformation \emph{déterministe} du nombre d'enregistrements en
montant, lisse mécaniquement la variance de sévérité. L'intervalle PRC est
donc plus resserré, non parce qu'il ignore l'incertitude de sévérité, mais
parce que cette incertitude y est intrinsèquement plus faible --- une nuance
reprise dans la décomposition de l'incertitude en Partie~5.

\vspace{1em}

% ==============================================================================
\section*{Comparaison méthodologique : limites des approches standards face au risque cyber}
\addcontentsline{toc}{section}{Comparaison méthodologique des approches}
% ==============================================================================

L'introduction d'un Modèle Interne Partiel (LDA) stochastique, tel que développé dans ce mémoire, représente un investissement méthodologique lourd pour un assureur. Pour en justifier la pertinence, il convient de le confronter aux deux méthodes traditionnellement utilisées sur la place pour évaluer le risque opérationnel : l'approche forfaitaire de la Formule Standard (Pilier 1) et l'approche déterministe par scénarios (historiquement privilégiée dans l'ORSA, Pilier 2). L'application de nos propres données de calibration à ces méthodes permet d'en démontrer les limites structurelles face à la nature du risque cyber.

%--------------------------------------------------------------
\subsection{La Formule Standard (Solvabilité II) : une mesure aveugle à la résilience}
%--------------------------------------------------------------

Dans le cadre de la Formule Standard de Solvabilité II (Article 204 du Règlement Délégué 2015/35), le risque cyber n'est pas isolé : il est noyé dans le module global du risque opérationnel ($SCR_{Op}$).

\paragraph{Rappel mathématique.} L'exigence de capital pour le risque opérationnel est calculée selon une approche modulaire, plafonnée à 30\,\% de l'exigence de capital de base ($BSCR$). En omettant la composante liée aux risques en unités de compte ($Exp_{ul}$), la formule se résume principalement à une fonction linéaire des primes acquises ($Earned_{premiums}$) et des provisions techniques ($TP$) :
\begin{equation}
SCR_{Op} \approx \min\Big( 0{,}3 \times BSCR \; ; \; \max(0{,}04 \times Earned_{premiums} \; ; \; 0{,}045 \times TP) \Big)
\end{equation}
(Note : la formule complète intègre également une pénalité marginale pour les croissances de primes annuelles supérieures à 20\,\%, que nous omettons ici par souci de clarté).

\paragraph{Démonstration numérique.} Appliquons cette formule à l'entité-type systémique utilisée pour le calibrage de nos résultats finaux (Partie 5), générant un volume de primes (Chiffre d'Affaires) de 15\,000 M\euro.
En appliquant le coefficient simplifié de 4\,\%, la charge de capital pour l'intégralité de ses risques opérationnels s'élève à :
\begin{equation}
SCR_{Op} = 0{,}04 \times 15\,000 = 600 \text{ M\euro}
\end{equation}
La confrontation avec notre modèle LDA-EVT est révélatrice. Sous la calibration PRC (plafonnée à 40~M\euro\ par sinistre, sévérité dérivée de la conversion Jacobs), la VaR 99,5\,\% pour une entité médiane est de \textbf{2\,772,7 M\euro}, soit déjà plus de quatre fois le forfait de Pilier~1 (600~M\euro). Face à la calibration OpRisk (capturant le risque extrême sans plafond, sur des montants réellement observés), la VaR atteint \textbf{9\,259,5 M\euro}. La Formule Standard sous-estime ainsi le risque de queue réel d'un facteur compris entre 4{,}6 (PRC) et 15 (OpRisk) pour cette entité systémique --- et c'est cette convergence des deux sources indépendantes vers un même verdict (un capital cyber en milliers de M\euro, très au-delà du Pilier~1), bien plus qu'un chiffre unique, qui constitue le message actuariel central du mémoire (cf.~Partie~5).

\paragraph{Synthèse comparative.} Le tableau ci-dessous résume cette confrontation pour notre entité de 15\,000 M\euro\ de CA, en observant le gain de capital obtenu en passant d'un profil non-conforme $S_2$ à un profil conforme $S_0$.

\begin{center}
\renewcommand{\arraystretch}{1.3}
\setlength{\tabcolsep}{4pt}
\begin{tabular}{>{\raggedright}p{4cm} >{\centering}p{2.5cm} >{\centering}p{3cm} >{\raggedright\arraybackslash}p{5cm}}
\toprule
\rowcolor{lightblue}
\textbf{Méthode d'évaluation} & \textbf{SCR Estimé (Médian)} & \textbf{Sensibilité DORA ($\Delta_{DORA}$)} & \textbf{Limites structurelles} \\
\midrule
Formule Standard (Pilier 1) & 600{,}0 M\euro & 0{,}0 M\euro & Aveugle à la résilience (aléa moral), sous-estimation extrême de la queue \\
\addlinespace
Déterministe (ORSA classique) & 200{,}0 M\euro & Difficile à quantifier & Biais d'ancrage expert, ignore l'agrégation et la contagion \\
\addlinespace
Modèle Interne Partiel (LDA -- PRC) & 2\,772{,}7 M\euro & 2\,014{,}7 M\euro & Nécessite un plafonnement opérationnel (espérance infinie) \\
\addlinespace
Modèle Interne Partiel (LDA -- OpRisk) & 9\,259{,}5 M\euro & 3\,879{,}3 M\euro & Biais de taille lié aux grandes entités financières historiques \\
\bottomrule
\end{tabular}
\end{center}

%--------------------------------------------------------------
\subsection{L'approche par scénarios déterministes (ORSA) : l'illusion de la précision}
%--------------------------------------------------------------

Pour pallier les lacunes de la Formule Standard, les assureurs s'appuient traditionnellement sur le Pilier 2 (ORSA) pour évaluer leurs risques extrêmes, souvent via une approche déterministe dite \og à dire d'expert\fg{}.

\paragraph{Mécanique de l'approche déterministe.} Le Directeur des Risques et le Responsable de la Sécurité des Systèmes d'Information (RSSI) définissent de concert un \og scénario catastrophe\fg{}. Par exemple : \emph{\og Une attaque par rançongiciel paralyse nos serveurs centraux pendant 5~jours, entraînant une perte d'exploitation, des coûts d'investigation et de notification.\fg{}} Les experts chiffrent ce scénario poste par poste, aboutissant à une perte ponctuelle, par exemple de 80 M\euro, qui est alors provisionnée comme exigence de capital.

\paragraph{Le biais de cognition face à la queue lourde.} Si cette méthode présente l'avantage d'être hautement \og lisible\fg{} pour la Direction Générale, nos données empiriques en démontrent la dangerosité statistique. L'expertise humaine a tendance à s'ancrer sur des moyennes ou des événements récents mémorisables. Dans la base SAS OpRisk Global, la moyenne des sinistres cyber est de 26,4 M\euro. Un expert fixant un scénario de crise à 80 M\euro\ (soit 3 fois la moyenne) pensera légitimement faire preuve d'un grand conservatisme.

Pourtant, l'analyse par la Théorie des Valeurs Extrêmes (EVT) menée en Partie 3 prouve le contraire. L'ajustement de la distribution de Pareto Généralisée (GPD) sur les données OpRisk a révélé un indice de queue $\hat{\xi} = 0{,}595$, caractéristique d'un régime à queue lourde extrême. Empiriquement, la base documente un sinistre maximum à \textbf{1\,500 M\euro}. En se contentant d'un scénario expert à 80 M\euro, l'assureur ignore la forme asymptotique de la distribution et ne provisionne absolument pas le quantile à 99,5\,\% exigé par la directive.

\paragraph{La supériorité de la simulation Monte-Carlo.} L'approche déterministe présente deux autres limites majeures que notre modèle LDA résout :
\begin{itemize}
    \item \textbf{La dépendance temporelle et l'agrégation :} Le scénario d'expert calcule le coût d'un incident isolé. Or, la VaR 99,5\,\% est une mesure annuelle. L'approche Monte-Carlo, en combinant des tirages de fréquence (loi binomiale négative) et de sévérité (GPD), permet d'explorer des années simulées où un sinistre majeur coïncide avec une myriade d'incidents d'attrition de taille moyenne, engendrant une perte agrégée supérieure à la simple addition déterministe.
    \item \textbf{La contagion externe (Tiers ICT) :} Un scénario construit en silo ignore la probabilité d'occurrence conjointe. L'introduction de la \textbf{copule de Gumbel} dans notre architecture permet de modéliser mathématiquement ce qu'un expert peine à quantifier : la survenance d'une crise systémique de marché ($\theta$ fortement négatif) qui déclenche à la fois la défaillance d'un fournisseur cloud critique et des failles internes. 
\end{itemize}

\paragraph{Conclusion sur les méthodes.} L'approche déterministe donne l'illusion de la précision en produisant un chiffre unique (ex: 80 M\euro), mais échoue à capturer le véritable risque de ruine. À l'inverse, l'approche stochastique développée dans ce mémoire assume l'incertitude (en produisant des intervalles de confiance par \textit{bootstrap}) et capture mathématiquement la réalité du risque cyber : une fréquence corrélée à l'environnement, une sévérité à queue lourde, et des pertes agrégées hautement non linéaires face à la qualité de la conformité DORA.

%==============================================================
\section{Résultats et discussion}
%==============================================================

Les quatre parties précédentes ont posé le cadre, les données, le moteur de
simulation et les scénarios de conformité. Cette cinquième partie présente
les résultats numériques obtenus, leur décomposition par source d'incertitude
et par brique de perte, et leur mise en perspective avec la littérature
existante et le pilotage prudentiel.

\begin{center}
\renewcommand{\arraystretch}{1.3}
\begin{tabular}{>{\raggedright}p{5.5cm} c >{\raggedright\arraybackslash}p{6cm}}
\toprule
\rowcolor{navy}
\textcolor{white}{\textbf{Paramètre}} & \textcolor{white}{\textbf{Valeur}} & \textcolor{white}{\textbf{Source / Justification}} \\
\midrule
\rowcolor{lightblue}\multicolumn{3}{l}{\textbf{Fréquence (Loi Binomiale Négative)}} \\
Fréquence de base ($\lambda_{ref}$) & 341 inc./an & PRC (Proxy sectoriel US) \\
Sur-dispersion ($\phi$) & 9{,}2 & Fixée pour maintenir la forme de variance \\
\addlinespace
\rowcolor{lightblue}\multicolumn{3}{l}{\textbf{Sévérité (GPD POT)}} \\
Seuil OpRisk ($u$) & 20 M\euro & Analyse MEP et stabilité $\xi$ \\
Seuil PRC ($u$) & 4{,}18 M\euro & Percentile 85 (sévérité dérivée Jacobs) \\
Indice de queue ($\hat{\xi}_{\text{OpRisk}}$) & 0{,}595 & MLE sur OpRisk (Moments finis) \\
Indice de queue ($\hat{\xi}_{\text{PRC}}$) & 1{,}03 & Conversion Jacobs (Espérance infinie) \\
Plafond opérationnel PRC & 40 M\euro & Capacité maximale réassurance \\
\addlinespace
\rowcolor{lightblue}\multicolumn{3}{l}{\textbf{Dépendance \& Variables latentes}} \\
Paramètre Gumbel ($\theta_{\text{Gumbel}}$) & 1{,}5 & Dépendance de queue supérieure modérée \\
Sensibilité systémique ($\rho_{\text{Vasicek}}$) & 0,2 à 0,8 & Calibré selon la criticité de l'entité \\
\bottomrule
\end{tabular}
\captionof{table}{Récapitulatif des paramètres centraux du modèle de simulation.}
\end{center}

%--------------------------------------------------------------
\subsection{Distribution du SCR\_DORA : ancrage sur une entité de 15 000 M\euro}
%--------------------------------------------------------------

Pour ancrer les quantiles simulés dans une réalité bilancielle et évaluer le risque de ruine, l'ensemble des modélisations a été exécuté sous l'hypothèse d'une entité de référence générant un Chiffre d'Affaires (CA) annuel de \textbf{15 000 M\euro} (assureur ou bancassureur de dimension européenne).

Le capital total à l'horizon 99{,}5~\% ($VaR_{99{,}5\%}$) et l'Expected Shortfall ($ES_{99{,}5\%}$) varient très fortement selon la source de sévérité retenue et le profil de conformité de l'entité.

\begin{center}
\renewcommand{\arraystretch}{1.3}
\setlength{\tabcolsep}{4pt}
\begin{tabular}{>{\raggedright}p{1.5cm} >{\raggedright}p{2.4cm} >{\centering}p{2.2cm} >{\centering}p{2.2cm} >{\centering}p{2cm} >{\centering\arraybackslash}p{2cm}}
\toprule
\rowcolor{lightblue}
\textbf{Source} & \textbf{Profil} & \textbf{VaR 99,5\,\%} & \textbf{ES 99,5\,\%} & \textbf{Ratio VaR/CA} & \textbf{Ratio ES/CA} \\
\midrule
PRC & Leader & 1\,778{,}7 M\euro & 1\,883{,}8 M\euro & 11{,}86\,\% & 12{,}56\,\% \\
PRC & Médian & 2\,772{,}7 M\euro & 2\,910{,}3 M\euro & 18{,}48\,\% & 19{,}40\,\% \\
PRC & Retardataire & 3\,724{,}3 M\euro & 3\,886{,}9 M\euro & 24{,}83\,\% & 25{,}91\,\% \\
\addlinespace
OpRisk & Leader & 6\,420{,}7 M\euro & 14\,255{,}0 M\euro & 42{,}80\,\% & 95{,}03\,\% \\
OpRisk & Médian & 9\,259{,}5 M\euro & 20\,023{,}2 M\euro & 61{,}73\,\% & 133{,}49\,\% \\
OpRisk & Retardataire & 10\,992{,}4 M\euro & 21\,720{,}1 M\euro & 73{,}28\,\% & 144{,}80\,\% \\
\bottomrule
\end{tabular}
\end{center}

Le SCR\_DORA se situe ainsi dans une fourchette globale de \textbf{1\,779~M\euro\ à près de 21\,720~M\euro} selon la source, le profil de conformité et la mesure de risque retenue.

\paragraph{Plausibilité économique et risque de ruine.} La lecture des ratios relatifs au Chiffre d'Affaires est particulièrement éloquente. Sous la calibration OpRisk, l'Expected Shortfall d'une entité non-conforme atteint \textbf{144,80\,\%} de son CA annuel ; même sous la calibration PRC, plus modérée, l'entité retardataire mobilise près de \textbf{26\,\%} de son CA en capital. En assurance classique, un sinistre dépassant le chiffre d'affaires est rare ; en risque cyber, une telle valeur traduit mathématiquement un risque de ruine certain de l'entité. Plutôt que d'invalider le modèle, ces montants confirment que la distribution sous-jacente (GPD à queue lourde, $\hat{\xi} \geq 0{,}6$ sous les deux sources) modélise un régime de pertes qui outrepasse la capacité d'absorption financière stricte d'une institution individuelle. C'est précisément ce constat actuariel qui justifie la philosophie du règlement DORA : imposer la résilience technique \textit{ex-ante} face à un risque dont l'impact \textit{ex-post} est financièrement intenable.

Cet écart, d'un facteur proche de~\textbf{12} entre la borne basse (PRC,
profil leader, VaR) et la borne haute (OpRisk, profil retardataire, ES),
demeure important, mais il s'est \emph{fortement resserré} par rapport à la
calibration PRC externe utilisée dans une version antérieure de ce travail
(où il atteignait un facteur~65). Ce resserrement est un résultat en soi :
une fois la PRC recalibrée de bout en bout sur ses données brutes (conversion
Jacobs sur 15\,053 incidents, cf.\ Partie~2), les deux sources
\emph{indépendantes} convergent sur le même diagnostic --- un capital cyber se
comptant en milliers de M\euro, très au-delà du forfait de Pilier~1. La
divergence résiduelle relève de l'incertitude légitime entre deux périmètres
distincts, non d'un désaccord de fond. Le SCR cyber ne peut donc toujours pas
être communiqué comme un point, mais la fourchette inter-sources est
désormais mieux maîtrisée.

Ce facteur~12 s'explique par le traitement de la queue de
distribution OpRisk : le paramètre de forme
$\hat{\xi}_{\text{OpRisk}} = 0{,}595 < 1$ garantissant une espérance finie
(cf.\ Partie~2), aucun plafonnement de la sévérité par sinistre n'y est
appliqué, à la différence de la PRC ($\hat{\xi}_{\text{PRC}} = 1{,}03 \geq 1$,
espérance infinie, où le plafond de 40~M\euro\ reste nécessaire à la
stabilité numérique de la simulation, cf.\ Partie~2.4). Le capital simulé
sous OpRisk reflète ainsi directement les événements extrêmes réellement
observés dans cette base (jusqu'à 1\,500~M\$ par sinistre), dont la
calibration GPD elle-même annonçait un $VaR_{99,5\%}$ \emph{par sinistre} de
662{,}8~M\euro --- déjà bien au-delà de tout plafond de 40~M\euro.

La figure~\ref{fig:distribution-scr} illustre la forme de la distribution
agrégée des pertes sous le régime OpRisk le plus sévère : une masse
concentrée aux faibles montants, prolongée d'une queue extrêmement lourde qui
repousse la VaR~99{,}5~\% très loin dans la distribution, matérialisant
graphiquement le poids désormais dominant des scénarios extrêmes dans la
constitution du capital.

\begin{figure}[htbp]
\centering
\includegraphics[width=0.9\textwidth]{figures/distribution_scr_oprisk_s2.png}
\caption{Distribution simulée des pertes agrégées (SCR cyber), source de
sévérité OpRisk, scénario de non-conformité totale $S_2$. La VaR~99{,}5~\%
(pointillés rouges) est repoussée très à droite de la masse principale,
signature de la queue lourde non plafonnée. La valeur exacte de la VaR
dépend du profil de conformité et du scénario retenus (cf.\ tableau
ci-dessus).}
\label{fig:distribution-scr}
\end{figure}



%--------------------------------------------------------------
\subsection{Décomposition de l'incertitude (variance par source)}
%--------------------------------------------------------------

Trois sources distinctes d'incertitude se superposent dans les résultats
présentés, et leur hiérarchie --- resserrée après la recalibration en propre
de la PRC --- est elle-même un résultat central du mémoire.

\paragraph{Le choix de source de sévérité reste le premier facteur, mais d'un
poids désormais comparable à celui de la conformité.} Entre PRC et OpRisk,
pour un même profil et une même mesure de risque, le capital total varie d'un
facteur proche de \textbf{3{,}3} (par exemple 2\,772{,}7~M\euro contre
9\,259{,}5~M\euro pour le profil médian en VaR). Cet écart demeure supérieur
au facteur associé au profil de conformité de l'entité (cf.\ paragraphe
suivant), mais il n'en est plus séparé d'un ordre de grandeur : c'est un
changement notable par rapport à une version antérieure de ce travail, où une
calibration PRC issue de paramètres externes non reproduits sur les données
brutes conduisait à un facteur inter-sources de~19. Une fois la PRC
recalibrée de bout en bout (conversion Jacobs $\ln(L)=7{,}68+0{,}76\ln(X)$ sur
les 15\,053 incidents 2019--2025, $\hat\xi_{\text{PRC}}=1{,}03$), l'écart de
niveau entre sources se réduit à un facteur~3, l'essentiel de l'ancien écart
provenant d'une sous-estimation de l'échelle de sévérité PRC plutôt que d'une
divergence de fond entre les deux bases. La différence résiduelle
s'explique par l'écart de régime de queue ($\hat{\xi}_{\text{PRC}} = 1{,}03$
avec plafond à 40~M\euro\ contre $\hat{\xi}_{\text{OpRisk}} = 0{,}595$ sans
plafond) et par le biais de taille d'OpRisk, qui documente des sinistres
extrêmes (jusqu'à 1\,500~M\$) que la sévérité PRC, plafonnée, ne peut
restituer.

\paragraph{Le profil de conformité de l'entité conserve un effet net, d'un
ordre désormais voisin de celui du choix de source.} Sous la source
OpRisk, le passage du profil leader au profil retardataire fait passer le
capital de 6\,420{,}7 à 10\,992{,}4~M\euro en VaR, soit un facteur~1{,}71 ;
sous PRC, le même passage fait passer le capital de 1\,778{,}7 à
3\,724{,}3~M\euro, soit un facteur~2{,}09. Ces facteurs, proches entre les deux
sources, confirment la robustesse \emph{qualitative} de l'effet de
conformité, indépendamment du niveau absolu de sévérité retenu : quelle que
soit la source de calibration, un profil de conformité dégradé multiplie le
capital requis par un facteur de l'ordre de~1{,}7 à~2{,}1. Cette stabilité du
facteur de conformité illustre que le modèle capture deux informations de
nature différente --- une sensibilité relative au profil, robuste au choix
de calibration, et un niveau absolu, qui l'est moins mais dont l'écart
inter-sources s'est lui aussi resserré.

\paragraph{L'environnement systémique amplifie la sensibilité au profil.}
La comparaison des scénarios $\theta = 0$ et $\theta = -2{,}5$ (choc cyber
sectoriel négatif) montre que la dégradation de l'environnement systémique
accroît le capital de toutes les entités, mais de façon différenciée : le
profil leader, sous OpRisk, voit son capital de base (6\,420{,}7~M\euro) s'accroître
sous l'effet du choc, alors que son $PCD$ passe de 0{,}0045 à 0{,}385 --- une
sensibilité au facteur commun directement héritée de la structure de Vasicek 
introduite en~1.4. Sous PRC, l'effet est de même ampleur. Le profil retardataire, 
déjà proche de $PCD = 1$ en régime normal, ne peut en revanche plus se dégrader 
significativement sous le choc (augmentation marginale sous OpRisk comme sous PRC), 
ce qui traduit un effet de saturation cohérent avec la forme de la fonction de lien $\Phi$.

La figure~\ref{fig:sensibilite-theta} rend visible ce mécanisme sur la
fréquence de sinistralité elle-même : à mesure que l'environnement se dégrade
($\theta$ décroissant), la fréquence de l'entité médiane croît fortement pour
rejoindre celle de l'entité en retard, tandis que le profil leader reste quasi
insensible au choc. C'est cette différenciation de sensibilité, et non un
simple décalage de niveau, qui constitue l'apport de la structure de Vasicek.

\begin{figure}[htbp]
\centering
\includegraphics[width=0.9\textwidth]{figures/sensibilite_theta.png}
\caption{Sensibilité de la fréquence de sinistralité $\lambda$ au choc
systémique $\theta$, par profil d'entité. Le profil leader (vert) reste quasi
plat ; le profil médian (orange) rattrape le profil en retard (rouge) sous un
choc sévère, illustrant l'effet de contagion différenciée hérité du modèle à
facteur unique. L'axe des abscisses va de $\theta = 0$ (environnement normal)
à $\theta = -3$ (crise majeure).}
\label{fig:sensibilite-theta}
\end{figure}

%--------------------------------------------------------------
\subsection{Allocation d'Euler : contribution de chaque brique}
%--------------------------------------------------------------

% NOTE(relecture) : deux corrections successives ont été apportées à la brique
% prestataire. (1) Le plafond de sévérité (40 M€) s'applique désormais à la
% fois à la remédiation ET au prestataire sous PRC (avant, seule la remédiation
% l'était sous PRC, xi>=1). (2) Cela ne suffisait pas : la brique prestataire
% était calibrée en Lognormale à échelle ABSOLUE (médiane ~5 M€), indépendante
% de la source active, alors que la remédiation change d'échelle selon la
% source (PRC/Jacobs vs OpRisk). Cette échelle absolue fixe écrasait la
% remédiation dès que les échelles différaient entre briques. La brique
% prestataire est désormais un SURCOÛT RELATIF appliqué à la sévérité de
% remédiation active (+8,3% à +11,8%, IBM/Ponemon Cost of a Data Breach 2025),
% ce qui restaure la commensurabilité entre briques et entre sources.
Cette décomposition révèle un enseignement qui n'apparaissait pas dans une
version antérieure du moteur, où la brique prestataire était calibrée sur une
échelle absolue (Lognormale, médiane $\approx 5$~M\euro) indépendante de la
source de sévérité active. Cette échelle absolue fixe entrait en conflit avec
la remédiation dès que la source active portait sa sévérité sur une échelle
différente, écrasant artificiellement la remédiation --- y compris
après plafonnement symétrique des deux briques à 40~M\euro, le problème
n'étant pas la queue mais le corps de la distribution. La brique prestataire
est donc désormais calibrée comme un \textbf{surcoût relatif} sur la sévérité
de remédiation active plutôt qu'une échelle absolue fixe : IBM/Ponemon
\emph{Cost of a Data Breach 2025} mesure un surcoût de~8{,}3~\% à~11{,}8~\%
lorsqu'un tiers est impliqué dans l'incident, fourchette tirée uniformément
comme les multiplicateurs DORA (Partie~4.3). Cette reformulation restaure la
commensurabilité entre briques \emph{et} entre sources : la brique
prestataire suit désormais l'échelle de la source active, au lieu d'importer
une échelle étrangère au régime simulé.

Sous les deux sources, la brique de \textbf{remédiation} domine nettement
(83{,}1 à 86{,}1~\% du capital selon le profil et la source), suivie par le
\textbf{prestataire} (13{,}9 à 16{,}9~\%). L'écart entre PRC et OpRisk sur
cette décomposition est désormais faible (moins de~3 points sur chaque
brique) : les deux sources, bien que très différentes en niveau absolu
(cf.~5.1--5.2), s'accordent maintenant sur la \emph{structure} du risque une
fois les briques rendues commensurables. C'est un résultat plus robuste, et
plus conforme à l'intuition actuarielle, que l'inversion PRC/OpRisk observée
avant cette correction.

La contribution de la \textbf{sanction} reste, sous les deux sources, très
en retrait (inférieure à~1~\% sous PRC, inférieure à~0{,}1~\% sous OpRisk),
en cohérence avec le poids structurellement borné du montant réglementaire
face aux deux autres briques.

La figure~\ref{fig:euler} synthétise visuellement cette hiérarchie. On
observe que le risque physique (remédiation) concentre l'essentiel du
capital sous les deux sources (83 à 86~\%), le risque tiers (prestataire)
en représentant une part minoritaire mais non négligeable (14 à 17~\%) ; la
sanction reste marginale sous les deux sources.

\begin{figure}[htbp]
\centering
\includegraphics[width=1\textwidth]{figures/euler_decomposition.png}
\caption{Décomposition du SCR (VaR~99{,}5~\%) par brique, sous les deux
sources de sévérité (profil médian, $\theta=0$). La brique de remédiation
domine sous les deux sources (85 à 86~\%) ; le prestataire, désormais
calibré en surcoût relatif sur la sévérité de remédiation active, représente
une part minoritaire et cohérente entre sources (14~\%).}
\label{fig:euler}
\end{figure}



\paragraph{Effet du plafond de sanction.} Une analyse de sensibilité
complémentaire confirme le caractère marginal de la brique de sanction dans
le capital total, y compris après correction du moteur. Le doublement du
plafond réglementaire de l'amende, de 20~M\euro à 40~M\euro, ne déplace que
très faiblement la VaR~99{,}5~\% : l'effet maximal observé reste inférieur à
7~M\euro\ (6{,}9~M\euro\ sous OpRisk et 4{,}4~M\euro\ sous PRC), y compris
pour le profil retardataire subissant un taux de non-conformité extrême ---
un effet d'autant plus négligeable en proportion que le capital total sous
OpRisk se compte désormais en milliers de millions d'euros.

\begin{center}
\renewcommand{\arraystretch}{1.3}
\begin{tabular}{>{\raggedright}p{2.2cm} >{\raggedright}p{2.5cm} >{\centering}p{3cm} >{\centering}p{3cm} >{\centering\arraybackslash}p{2.6cm}}
\toprule
\rowcolor{lightblue}
\textbf{Source} & \textbf{Profil} & \textbf{VaR base (plafond 20~M\euro)} & \textbf{VaR stress (plafond 40~M\euro)} & \textbf{$\Delta$VaR} \\
\midrule
OpRisk & Leader & 6\,420{,}7 M\euro & 6\,420{,}7 M\euro & 0{,}0 M\euro \\
OpRisk & Médian & 9\,259{,}5 M\euro & 9\,260{,}3 M\euro & 0{,}9 M\euro \\
OpRisk & Retardataire & 10\,992{,}4 M\euro & 10\,999{,}3 M\euro & 6{,}9 M\euro \\
\addlinespace
PRC & Leader & 1\,778{,}7 M\euro & 1\,778{,}7 M\euro & 0{,}0 M\euro \\
PRC & Médian & 2\,772{,}7 M\euro & 2\,775{,}5 M\euro & 2{,}9 M\euro \\
PRC & Retardataire & 3\,724{,}3 M\euro & 3\,728{,}7 M\euro & 4{,}4 M\euro \\
\bottomrule
\end{tabular}
\end{center}

Cette insensibilité confirme un message central du mémoire pour le pilotage :
le capital cyber répond à un risque \emph{opérationnel} (fréquence et
sévérité des incidents), et non à un risque de \emph{sanction} administrative,
dont l'impact financier demeure borné par sa structure réglementaire et
négligeable face à la queue lourde des pertes techniques, sous les deux
sources.

%--------------------------------------------------------------
\subsection{Analyse de sensibilité aux hypothèses porteuses}
%--------------------------------------------------------------

Le modèle mobilise des paramètres de statut épistémique hétérogène : certains
sont \emph{calibrés sur données} et assortis d'un intervalle de confiance
(l'indice de queue $\hat\xi$ et l'échelle $\hat\sigma$ de la GPD, dont
l'IC90\% est obtenu par bootstrap, cf.\ Partie~2) ; d'autres sont des
\emph{hypothèses libres}, ancrées sur proxies ou avis d'expert et non estimées
faute de données conditionnelles à la conformité (la corrélation systémique
$\gamma$ du modèle de Vasicek, la non-conformité de référence $p_{ref}$, les
bornes des multiplicateurs de fréquence, la surcharge tiers, la sur-dispersion
$\varphi$, le paramètre de la copule). La légitimité du modèle exige de
mesurer \emph{combien} chacun de ces leviers pèse réellement sur le capital,
plutôt que de se contenter de les déclarer incertains.

Nous perturbons donc chaque paramètre sur une plage plausible et documentée,
toutes choses égales par ailleurs, et mesurons la réponse de la
VaR~99{,}5~\% autour du cas de référence (OpRisk, profil médian, $\theta=0$,
soit un SCR de base de 9\,259~M\euro). La figure~\ref{fig:tornado} classe les
paramètres par amplitude d'effet (\emph{tornado}).

\begin{figure}[htbp]
\centering
\includegraphics[width=\textwidth]{figures/tornado_sensitivity.png}
\caption{Sensibilité de la VaR~99{,}5~\% aux hypothèses porteuses du modèle,
cas de référence OpRisk / profil médian / $\theta=0$. En bleu, les paramètres
calibrés sur données (IC90\% bootstrap) ; en orange, les hypothèses libres.
Le paramètre $\gamma$ étant inerte à $\theta=0$ par construction (terme
$\gamma\cdot\theta$), sa sensibilité est évaluée sous choc systémique
$\theta=-2{,}5$, d'où le décalage de sa barre.}
\label{fig:tornado}
\end{figure}

\paragraph{Un résultat qui conforte la thèse du mémoire.} La hiérarchie est
sans appel et va dans le sens de tout le reste du travail :
\begin{itemize}[topsep=3pt,itemsep=1pt]
\item l'indice de queue $\hat\xi$ domine \emph{massivement} : sur son seul
IC90\% bootstrap $[0{,}30\,;\,0{,}83]$, la VaR passe de 3\,073 à
32\,312~M\euro, soit une amplitude de plus de 300\,\% du SCR de base --- à lui
seul, le paramètre de queue peut multiplier le capital par dix ;
\item l'échelle $\hat\sigma$ vient ensuite ($\pm 69\,\%$) ; les deux
paramètres calibrés sur données concentrent ainsi l'essentiel de
l'incertitude ;
\item parmi les hypothèses libres, seule la non-conformité de référence
$p_{ref}$ a un effet notable ($\pm 18\,\%$) ; toutes les autres
--- multiplicateurs de fréquence ($\pm 5\,\%$), sur-dispersion, corrélation
systémique $\gamma$ même sous choc ($\sim 3\,\%$), surcharge tiers et copule
($< 1\,\%$) --- déplacent le capital de façon marginale.
\end{itemize}

Cette analyse préempte directement la critique naturelle adressée à un modèle
partiellement calibré à dire d'expert : les paramètres les moins solidement
étayés ($\gamma$, la copule, la surcharge tiers) sont précisément ceux qui
pèsent le moins sur le résultat, tandis que le déterminant premier du capital
--- l'indice de queue de la sévérité --- est estimé sur données réelles et
assorti de son intervalle de confiance. Le corollaire rejoint le message
central des sections~5.1--5.2 : l'incertitude du SCR cyber est d'abord une
incertitude \emph{de sévérité} (donc de source de données), très loin devant
l'incertitude sur la mécanique de conformité. Investir dans une meilleure
donnée de sévérité réduirait l'incertitude du capital bien davantage que
raffiner la calibration du pont latent.

\paragraph{Validation et perspectives de robustesse.} L'absence de données de
pertes conditionnelles à la conformité DORA interdit, à ce stade, une
validation \emph{out-of-sample} au sens strict : on ne peut confronter le
$\Delta_{DORA}$ prédit à un $\Delta_{DORA}$ réalisé, faute d'historique. Trois
formes de validation restent néanmoins accessibles et sont mobilisées dans ce
travail. D'abord, une validation \emph{croisée entre sources} : la calibration
de queue est confrontée à une source indépendante (PRC vs OpRisk), et leur
écart --- plutôt que masqué --- est traité comme une composante de
l'incertitude, ce qui constitue un test de robustesse du régime de queue
(toutes deux confirment $\xi > 0$, cf.\ Partie~2). Ensuite, une validation
\emph{de cohérence interne} : la décomposition d'Euler, une fois les briques
rendues commensurables, converge entre sources sur la structure du risque
(remédiation dominante, cf.\ 5.3), signe que le résultat n'est pas un artefact
d'une source particulière. Enfin, l'analyse de sensibilité ci-dessus joue le
rôle d'un \emph{stress-test de spécification} : elle borne l'effet de chaque
hypothèse et localise l'incertitude résiduelle. À terme, la constitution du
registre d'incidents DORA (cf.\ Partie~2.6) ouvrira la voie à une validation
directe : rétro-tester la calibration de sévérité sur les pertes nouvellement
observées, et surtout estimer les multiplicateurs de fréquence sur des comptes
de sinistres réellement conditionnels au niveau de conformité --- remplaçant
alors l'hypothèse experte porteuse par une quantité mesurée. Le modèle est
construit pour absorber cette substitution sans changement d'architecture.

%--------------------------------------------------------------
\subsection{Robustesse de la calibration de sévérité : validation temporelle et adéquation formelle}
%--------------------------------------------------------------

Les paragraphes précédents évoquaient une validation \emph{de cohérence
interne} et un \emph{stress-test de spécification}, sans toutefois soumettre
la calibration GPD elle-même à un test statistique formel ni à une validation
hors échantillon. Les deux diagnostics suivants comblent ce trou, en
n'utilisant que les incidents OpRisk déjà mobilisés --- aucune donnée
supplémentaire.

\paragraph{Validation par découpage temporel.} La base OpRisk cyber$\times$finance
(570 incidents, 2000--2026) est scindée à la date de coupure 2017
(349 incidents avant, 221 après), et la GPD est recalibrée indépendamment sur
chaque période, seuil de percentile~85 propre à chacune :

\begin{center}
\renewcommand{\arraystretch}{1.25}
\begin{tabular}{l r r r r}
\toprule
\rowcolor{lightblue}
\textbf{Période} & \textbf{Seuil $u$} & \textbf{$n$ excès} & \textbf{$\hat{\xi}$ (IC90\%)} & \textbf{$\hat{\sigma}$} \\
\midrule
Entraînement ($\leq 2017$) & 11{,}27~M\euro & 53 & $0{,}71\ [0{,}25\,;\,1{,}17]$ & 25{,}10~M\euro \\
Test ($> 2017$)            & 49{,}38~M\euro & 33 & $0{,}57\ [0{,}12\,;\,1{,}01]$ & 103{,}63~M\euro \\
\bottomrule
\end{tabular}
\end{center}

Les intervalles de confiance sur $\xi$ se chevauchent largement : la
\emph{forme} de la queue (régime Fréchet, $\xi>0$) est stable entre les deux
périodes, ce qui valide le choix structurel de la GPD lui-même. En revanche,
l'\emph{échelle} $\hat{\sigma}$ quadruple d'une période à l'autre (25,10 à
103,63~M\euro), et un test de Kolmogorov-Smirnov --- ici pleinement valide,
les paramètres du modèle testé provenant exclusivement de la période
d'entraînement, indépendante de l'échantillon testé --- confirme que les
excès de la période de test au-delà du seuil d'entraînement ($n=69$,
statistique KS$=0{,}168$, $p=0{,}037$) ne suivent pas la loi
$GPD(\hat{\xi}_{\text{train}},\hat{\sigma}_{\text{train}})$ : la VaR~$99{,}5\,\%$
extrapolée à partir du seul régime 2000--2017 (371,1~M\euro) sous-estime
largement le quantile~$99{,}5\,\%$ effectivement observé sur 2018--2026
(1\,077,5~M\euro). Ce résultat ne fragilise pas le modèle : il \emph{confirme
quantitativement}, par un test formel, ce que la section~2.3 affirmait déjà
qualitativement sur l'inflation cyber --- calibrer sur des données trop
anciennes sous-estime la sévérité actuelle, non parce que le régime de queue
change de nature, mais parce que son échelle croît dans le temps. La
calibration officielle du mémoire (section~2.7), qui retient l'intégralité de
la période 2000--2026 plutôt qu'une sous-période ancienne, se trouve ainsi
justifiée \emph{a posteriori} par ce test, plutôt que par la seule
plausibilité qualitative.

\begin{figure}[htbp]
\centering
\includegraphics[width=0.7\textwidth]{figures/holdout_validation.png}
\caption{Validation holdout : fonction de survie empirique des excès de la
période de test (au-delà du seuil d'entraînement) comparée à la prédiction
$GPD(\hat{\xi}_{\text{train}}, \hat{\sigma}_{\text{train}})$. L'écart visible
en échelle log matérialise le sous-provisionnement qu'aurait produit une
calibration figée sur la seule période 2000--2017.}
\label{fig:holdout-validation}
\end{figure}

\paragraph{Adéquation formelle de la calibration officielle.} Sur
l'intégralité du périmètre (570 incidents, seuil $u=20{,}03$~M\euro, 86
excès --- reproduisant à moins de $0{,}1\,\%$ près les $\hat{\xi}=0{,}595$ et
$\hat{\sigma}=57{,}97$~M\euro déjà publiés en~2.7), la qualité de
l'ajustement GPD, jusqu'ici appuyée sur une lecture graphique (mean excess
plot, diagnostic de stabilité, Hill plot), est soumise à deux tests
quantitatifs. La transformée intégrale de probabilité des excès --- qui doit
suivre une loi Uniforme$(0,1)$ sous un modèle bien spécifié --- donne une
statistique de Kolmogorov-Smirnov de $0{,}046$ ($p\approx0{,}99$) et une
statistique d'Anderson-Darling de $0{,}14$, très en-deçà du repère critique
usuel à 5\,\% ($\approx 2{,}49$) : l'hypothèse d'adéquation n'est pas rejetée.
Cette lecture appelle une réserve méthodologique explicite : les paramètres
$\hat{\xi}$ et $\hat{\sigma}$ étant estimés sur le \emph{même} échantillon que
celui soumis au test, la loi exacte de la statistique KS diffère légèrement
du cas i.i.d.\ à paramètres connus a priori --- les $p$-valeurs rapportées
sont donc indicatives, non un test exact au sens strict. Elles corroborent
néanmoins, de façon convergente, le diagnostic graphique déjà présenté en~2.7
et~3.3 : le seuil retenu (percentile~85) ne se contente pas de stabiliser
visuellement $\hat{\xi}$, il produit un ajustement statistiquement défendable.

\begin{figure}[htbp]
\centering
\includegraphics[width=0.95\textwidth]{figures/gpd_goodness_of_fit.png}
\caption{Adéquation de la calibration GPD officielle (OpRisk, $u=20{,}03$~M\euro).
\emph{Gauche} : QQ-plot des excès contre les quantiles théoriques de la GPD
ajustée. \emph{Droite} : histogramme de la transformée intégrale de
probabilité, comparé à la densité Uniforme$(0,1)$ attendue sous adéquation.}
\label{fig:gpd-goodness-of-fit}
\end{figure}

%--------------------------------------------------------------
\subsection{Robustesse du moteur de simulation : famille de copule et convergence Monte Carlo}
%--------------------------------------------------------------

Les deux diagnostics précédents bornent l'incertitude de \emph{paramètre}.
Restent deux questions de nature différente, jusqu'ici non instruites dans ce
mémoire : le choix de la \emph{forme fonctionnelle} de la copule pèse-t-il sur
le résultat autant que son intensité de dépendance~? Et le nombre de tirages
Monte Carlo retenu dans l'ensemble du pipeline ($n_{sim}=100\,000$) suffit-il
à garantir que le SCR rapporté ne soit pas lui-même bruité par l'erreur
d'échantillonnage de la simulation~? Ces deux tests ne mobilisent aucune
donnée supplémentaire : ils exploitent uniquement le moteur de simulation et
les calibrations déjà fixées en Parties~2--3.

\paragraph{Robustesse à la famille de copule.} La copule de Gumbel est
justifiée en~1.9 par sa dépendance de queue supérieure, sans qu'aucun test
quantitatif ne mesure si le choix de cette \emph{famille} pèse sur le capital
autant que son \emph{paramètre} (déjà testé en~5.4, effet $<1\,\%$). Trois
familles alternatives sont ici confrontées à Gumbel, chacune appariée sur le
\emph{même} tau de Kendall en régime non conforme ($\tau=0{,}444$, soit
$\theta_{\text{Gumbel}}=1{,}8$) et conforme ($\tau=0{,}167$,
$\theta_{\text{Gumbel}}=1{,}2$), de sorte que seule la forme fonctionnelle de
la dépendance varie, à intensité de dépendance strictement égale : Clayton
roté à $180°$ (dépendance de queue supérieure, forme différente de Gumbel),
Frank (\emph{aucune} dépendance de queue --- cas de contraste), et Student
($df=4$, dépendance de queue symétrique haute/basse). Le tableau ci-dessous
compare, pour chaque source et chaque famille, la VaR~$99{,}5\,\%$ en
scénario~S2 (non-conformité totale, profil médian) et le $\Delta_{DORA}$
associé, obtenus par le même moteur de simulation que le reste du mémoire ---
seule la famille de copule change d'une ligne à l'autre.

\begin{center}
\renewcommand{\arraystretch}{1.25}
\begin{tabular}{l r r r r}
\toprule
\rowcolor{lightblue}
\textbf{Famille (PRC)} & \textbf{VaR\,99,5\,\%} & \textbf{$\Delta_{DORA}$} & \textbf{Écart VaR} & \textbf{Écart $\Delta_{DORA}$} \\
\midrule
Gumbel (référence)     & 3\,741{,}9~M\euro  & 1\,960{,}4~M\euro & --- & --- \\
Clayton (roté)          & 3\,750{,}4~M\euro  & 1\,968{,}8~M\euro & $+0{,}2\,\%$ & $+0{,}4\,\%$ \\
Frank                   & 3\,752{,}6~M\euro  & 1\,974{,}6~M\euro & $+0{,}3\,\%$ & $+0{,}7\,\%$ \\
Student ($df=4$)        & 3\,745{,}4~M\euro  & 1\,967{,}1~M\euro & $+0{,}1\,\%$ & $+0{,}3\,\%$ \\
\midrule
\rowcolor{lightblue}
\textbf{Famille (OpRisk)} & \textbf{VaR\,99,5\,\%} & \textbf{$\Delta_{DORA}$} & \textbf{Écart VaR} & \textbf{Écart $\Delta_{DORA}$} \\
\midrule
Gumbel (référence)     & 10\,903{,}2~M\euro & 4\,602{,}3~M\euro & --- & --- \\
Clayton (roté)          & 10\,792{,}8~M\euro & 4\,523{,}2~M\euro & $-1{,}0\,\%$ & $-1{,}7\,\%$ \\
Frank                   & 10\,879{,}2~M\euro & 4\,543{,}8~M\euro & $-0{,}2\,\%$ & $-1{,}3\,\%$ \\
Student ($df=4$)        & 10\,877{,}4~M\euro & 4\,578{,}6~M\euro & $-0{,}2\,\%$ & $-0{,}5\,\%$ \\
\bottomrule
\end{tabular}
\end{center}

Ces montants sont des points de simulation directe de l'engin (scénario~S2,
profil médian), non passés par le bootstrap à deux niveaux de la Partie~4 :
ils diffèrent donc légèrement des valeurs bootstrappées déjà publiées en~4.4
et~5.1, et ne s'y substituent pas. Ce qui importe ici est l'\emph{écart
relatif} entre familles, à tau de Kendall égal : il reste, sous les deux
sources, inférieur à $1{,}7\,\%$ en valeur absolue --- du même ordre de
grandeur que l'effet du \emph{paramètre} de la copule déjà mesuré en~5.4
($<1\,\%$). Y compris face à Frank, qui ne comporte structurellement aucune
dépendance de queue, l'écart reste marginal : c'est l'\emph{intensité} de la
dépendance (le tau de Kendall, hérité de la contagion documentée en~1.5) qui
gouverne le capital, pas la forme fonctionnelle retenue pour la représenter.
Le choix de Gumbel, motivé qualitativement en~1.9, se trouve ainsi confirmé
quantitativement comme non déterminant pour le niveau de capital rapporté ---
un résultat rassurant pour la robustesse de l'architecture, obtenu sans
aucune donnée supplémentaire.

\begin{figure}[htbp]
\centering
\includegraphics[width=0.95\textwidth]{figures/copula_robustness.png}
\caption{VaR~$99,5\,\%$ (scénario S2, profil médian) sous quatre familles de
copule appariées sur le même tau de Kendall que Gumbel, par source de
sévérité. La ligne pointillée marque la référence Gumbel : l'écart entre
familles reste, sous les deux sources, inférieur à $1{,}7\,\%$.}
\label{fig:copula-robustness}
\end{figure}

\paragraph{Convergence de la simulation Monte Carlo.} L'ensemble du pipeline
retient $n_{sim}=100\,000$ années simulées sans qu'aucune vérification
n'établisse que ce choix suffit à stabiliser la VaR~$99{,}5\,\%$ rapportée.
Pour le mesurer, on répète l'estimation --- à paramètres strictement figés,
scénario~S2, profil médian --- sous plusieurs graines indépendantes, pour une
grille croissante de $n_{sim}$, et l'on observe la dispersion résiduelle
(l'erreur Monte Carlo pure, distincte de l'incertitude de paramètre déjà
traitée en~5.2 et~5.4) :

\begin{center}
\renewcommand{\arraystretch}{1.25}
\begin{tabular}{r r r r}
\toprule
\rowcolor{lightblue}
\textbf{$n_{sim}$} & \textbf{VaR moy.\ (PRC)} & \textbf{CV (PRC)} & \textbf{CV (OpRisk)} \\
\midrule
1\,000    & 3\,729{,}9~M\euro & $2{,}28\,\%$  & $20{,}83\,\%$ \\
3\,000    & 3\,744{,}5~M\euro & $0{,}98\,\%$  & $12{,}90\,\%$ \\
10\,000   & 3\,744{,}6~M\euro & $0{,}60\,\%$  & $8{,}03\,\%$ \\
30\,000   & 3\,743{,}5~M\euro & $0{,}36\,\%$  & $3{,}51\,\%$ \\
\textbf{100\,000}  & \textbf{3\,743{,}6~M\euro} & \textbf{0{,}30\,\%} & \textbf{1{,}38\,\%} \\
300\,000  & 3\,744{,}0~M\euro & $0{,}08\,\%$  & $1{,}15\,\%$ \\
\bottomrule
\end{tabular}
\end{center}

Le coefficient de variation Monte Carlo (CV, écart-type inter-graines rapporté
à la VaR moyenne) décroît de façon monotone avec $n_{sim}$, comme attendu pour
un estimateur convergent. À $n_{sim}=100\,000$ --- la valeur retenue dans
l'ensemble des notebooks du pipeline (\texttt{04} à \texttt{08}) --- le CV
s'établit à $0{,}23\,\%$ sous PRC et $1{,}38\,\%$ sous OpRisk, cette dernière
plus élevée du fait de la queue plus lourde de la source. Dans les deux cas,
cette erreur d'échantillonnage reste d'un ordre de grandeur très inférieur à
l'incertitude de paramètre documentée en~5.4 (amplitude bootstrap sur $\xi$
seul : plus de $300\,\%$ du SCR de base). Ce résultat valide \emph{a
posteriori} le choix de $n_{sim}=100\,000$ retenu dans tout le pipeline : au-delà
de ce seuil (300\,000), le CV se stabilise sans gain perceptible, signe que
l'effort de calcul supplémentaire ne réduirait pas l'incertitude du SCR ---
qui reste, comme conclu en~5.2 et~5.4, une incertitude de \emph{modèle et de
donnée}, non de simulation.

\begin{figure}[htbp]
\centering
\includegraphics[width=0.95\textwidth]{figures/mc_convergence.png}
\caption{Convergence de la VaR~$99,5\,\%$ (scénario S2, profil médian) avec le
nombre d'années simulées $n_{sim}$, par source. La bande ombrée représente
l'étendue (min--max) observée sur les graines indépendantes à chaque taille
d'échantillon ; la ligne pointillée marque la référence à grande taille
($n_{sim}=300\,000$).}
\label{fig:mc-convergence}
\end{figure}

%--------------------------------------------------------------
\subsection{Comparaison avec les travaux antérieurs}
%--------------------------------------------------------------

Les résultats obtenus s'inscrivent en cohérence, mais aussi en tension
partielle, avec les travaux antérieurs sur la structure du risque cyber
opérationnel. Le tableau ci-dessous confronte l'indice de queue $\hat\xi$
obtenu dans ce mémoire à ceux de la littérature actuarielle récente, le
paramètre de forme constituant le déterminant premier du niveau de capital.

\begin{center}
\renewcommand{\arraystretch}{1.3}
\begin{tabular}{>{\raggedright}p{5.5cm} >{\centering}p{3.2cm} >{\centering\arraybackslash}p{3.2cm}}
\toprule
\rowcolor{lightblue}
\textbf{Étude / source} & \textbf{$\hat\xi$ (indice de queue)} & \textbf{Espérance} \\
\midrule
Présent mémoire --- PRC (via Jacobs) & 1{,}03 \; [0{,}97 ; 1{,}09] & Infinie ($\xi \geq 1$) \\
Présent mémoire --- OpRisk Global & 0{,}595 \; [0{,}31 ; 0{,}83] & Finie \\
\addlinespace
Dacorogna, Debbabi \& Kratz (2023) --- base SCRC & 0{,}983 \; [0{,}93 ; 1{,}04] & Frontière $\xi \approx 1$ \\
Farkas, Lopez \& Thomas (2021) & \multicolumn{2}{c}{Espérance conditionnelle infinie sur certains profils} \\
Malavasi et al.\ (2022) & \multicolumn{2}{c}{Espérance infinie confirmée sur données réelles} \\
\bottomrule
\end{tabular}
\end{center}

\paragraph{Un régime de queue confirmé, une ampleur qui reste débattue.}
Deux enseignements se dégagent de cette confrontation. D'une part, la
calibration OpRisk ($\hat\xi \approx 0{,}6$) se situe dans les ordres de
grandeur des études de risque opérationnel bancaire sous approche AMA,
cohérente avec un régime de queue lourde mais à moments finis. D'autre part,
et c'est le point le plus significatif, la calibration PRC ($\hat\xi = 1{,}03$,
espérance infinie) n'est nullement une anomalie isolée : elle rejoint un
courant robuste de la littérature qui documente une \emph{espérance de
sévérité infinie} sur données réelles. Le rapprochement avec Dacorogna,
Debbabi et Kratz (2023) est ici frappant : sur la base des plaintes de la
Gendarmerie nationale, ces auteurs estiment un $\hat\xi$ à la frontière de
l'unité ($0{,}983$, IC recouvrant $1$) --- soit une valeur statistiquement
indiscernable de notre estimation PRC ($1{,}03$, IC $[0{,}97~;~1{,}09]$),
alors que les deux calibrations reposent sur des données, des juridictions et
des méthodes entièrement distinctes. Farkas, Lopez et Thomas (2021)
obtiennent de même, par arbres de régression Pareto généralisés, des
espérances conditionnelles infinies sur certains profils de sinistres, et
Malavasi et al.\ (2022) confirment ce résultat sur un jeu de données
indépendant.

\paragraph{Positionnement du mémoire.} Loin de trancher artificiellement ce
débat, la présente contribution l'\emph{encadre} : nos deux sources bornent
précisément l'intervalle des régimes de queue rencontrés dans la littérature,
la PRC matérialisant le régime à espérance infinie (protection des données)
et OpRisk le régime à moments finis (risque opérationnel bancaire). Cette
tension entre sources n'invalide pas les travaux antérieurs pris séparément ;
elle illustre au contraire que la littérature sur le risque cyber demeure
structurellement fragmentée selon le type de base mobilisée. En confrontant
explicitement deux sources plutôt qu'en n'en retenant qu'une, et en traitant
leur écart comme une composante à part entière de l'incertitude --- une
composante qui, après recalibration en propre de la PRC, se révèle du même
ordre que l'écart entre scénarios de conformité (facteur~1{,}9 entre sources
contre~2{,}6 entre scénarios sous OpRisk, cf.~Partie~5.1--5.2), et non plus
d'un ordre de grandeur au-dessus ---, le présent mémoire documente cette
fragmentation plutôt que de la masquer.

%--------------------------------------------------------------
\subsection{Implications pour l'ORSA et le pilotage réglementaire}
%--------------------------------------------------------------

Les résultats obtenus ont quatre implications directes pour le pilotage
prudentiel en Pilier~2.

Premièrement, la nette dominance de la brique de remédiation sous
\emph{les deux} sources (85 à 86~\% du capital, cf.~5.3) suggère que les
efforts de maîtrise du risque cyber devraient prioriser la réduction de la
fréquence et de la sévérité des incidents opérationnels majeurs, avant la
gouvernance des tiers ICT (13 à 17~\%) ou le risque de sanction (marginal).
Fait notable, cette hiérarchie des priorités est désormais \emph{robuste au
choix de source} : une fois la brique prestataire recalibrée en surcoût
relatif et la PRC recalibrée en propre, PRC et OpRisk s'accordent sur la
même structure à moins de trois points près (cf.~5.3), là où une version
antérieure du moteur produisait une inversion artefactuelle entre sources.
La convergence structurelle des deux calibrations indépendantes renforce donc
la confiance dans cette hiérarchisation des priorités de maîtrise.

Deuxièmement, la sensibilité du capital au facteur systémique $\theta$
documentée en~5.2 justifie l'intégration, dans les scénarios de stress de
l'ORSA, d'un scénario de dégradation collective de l'environnement cyber,
distinct des scénarios de dégradation individuelle de la gouvernance : les
deux mécanismes n'affectent pas le capital de la même façon, ni avec la
même intensité selon le profil de l'entité.

Troisièmement, l'écart entre sources de sévérité --- un facteur de l'ordre
de \textbf{3} sur le niveau de capital (profil médian, VaR), du même ordre
que le facteur~1{,}7 à~2{,}1 associé au profil de conformité --- plaide
pour que le pilotage interne ne retienne pas un unique scénario central, mais
documente explicitement, dans le rapport ORSA, la sensibilité du SCR cyber
au choix méthodologique de calibration. Cette pratique de transparence sur
l'incertitude de modèle, que la formule standard de Solvabilité~II,
forfaitaire, ne permet structurellement pas d'exprimer, apparaît comme
une implication significative du mémoire pour la gouvernance
du risque cyber en Pilier~2 : le choix de la base de calibration n'est pas
un détail technique, mais l'un des deux déterminants de premier ordre du
niveau de capital, à parité désormais avec le profil de conformité.

Quatrièmement, la marginalité confirmée de la brique de sanction --- sous
les deux sources, et de façon encore plus nette sous OpRisk --- conforte le
message selon lequel le capital cyber répond avant tout à un risque
opérationnel et non à un risque réglementaire ponctuel, dont l'ampleur reste
structurellement bornée par les plafonds légaux.

%==============================================================
\section*{Conclusion}
\addcontentsline{toc}{section}{Conclusion}
%==============================================================

Le SCR cyber n'est pas un nombre, c'est une distribution large. Les
résultats de ce mémoire le confirment à chaque étape : une fourchette de
capital total s'étendant de \textbf{1\,779~M\euro à près de 21\,720~M\euro} selon la
source de sévérité, le profil de conformité et la mesure de risque retenue,
soit un facteur de \textbf{12} entre ses bornes ; un $\Delta_{DORA}$ dont
l'intervalle de confiance s'élargit massivement sous la source
OpRisk ; et une allocation d'Euler qui, une fois la brique prestataire
rendue commensurable avec la remédiation (surcoût relatif plutôt qu'échelle
absolue, cf.~5.3), converge vers une structure cohérente entre sources : la
brique de remédiation porte l'essentiel du capital (83 à 86~\%) sous PRC
comme sous OpRisk. Une fois la PRC recalibrée de bout en bout sur ses
données brutes (conversion Jacobs, $\hat\xi_{\text{PRC}} = 1{,}03$), le
désaccord entre sources s'est en outre nettement resserré --- d'un facteur~65
à un facteur~12 entre bornes extrêmes, et d'un facteur~19 à un facteur~3 à
profil comparable : deux sources indépendantes qui, calibrées chacune avec
rigueur, convergent bien davantage qu'attendu vers un même diagnostic. Le
désaccord résiduel porte sur le \emph{niveau} de capital, non sur sa
\emph{structure} ni sur l'ordre de grandeur --- une information honnête sur
l'état actuel des données disponibles pour quantifier le risque cyber
réglementaire, plutôt qu'un échec de la modélisation.

Les perspectives de ce travail sont directement liées à cette limite. Le
registre d'incidents DORA que les autorités européennes de surveillance
(ESA) doivent progressivement constituer à partir des notifications
obligatoires prévues par le règlement constituera, à terme, la première
base réellement harmonisée à l'échelle européenne pour calibrer la
fréquence et la sévérité du risque cyber réglementaire --- une base qui
fait aujourd'hui structurellement défaut, et dont l'absence explique une
part substantielle de la largeur des intervalles présentés dans ce mémoire.
En ce sens, la contribution de ce travail n'est pas de clore le débat sur le
niveau du SCR cyber, mais d'en formaliser la structure d'incertitude,
dans l'attente de données futures plus riches.


\vspace{1em}


\bibliographystyle{apalike}
\addcontentsline{toc}{section}{Références}
\bibliography{references}
\end{document}